%% paper/main.tex
%% Main paper file for:
%%   "Revisiting the Compact Tail Assignment Model:
%%    Feasibility and Type-A Maintenance"
%%
%% Compile with:  latexmk -pdf main.tex
%% or:            pdflatex main && bibtex main && pdflatex main && pdflatex main
%% ==========================================================================
\documentclass[review, 12pt, authoryear]{elsarticle}

\usepackage[utf8]{inputenc}
\usepackage[T1]{fontenc}
\usepackage{lmodern}
\usepackage{amsmath, amssymb, amsthm, mathtools}
\usepackage{booktabs}
\usepackage{multirow}
\usepackage{tabularx}
\usepackage{graphicx}
\usepackage[ruled, vlined, linesnumbered]{algorithm2e}
\usepackage{hyperref}
\usepackage{cleveref}
\usepackage{xcolor}
\usepackage{siunitx}    % consistent number formatting in tables
\usepackage{caption}
\usepackage{subcaption}

%% theorem environments
\theoremstyle{plain}
\newtheorem{proposition}{Proposition}
\newtheorem{lemma}{Lemma}
\newtheorem{theorem}{Theorem}
\theoremstyle{definition}
\newtheorem{definition}{Definition}
\theoremstyle{remark}
\newtheorem{remark}{Remark}

\newtheorem{example}{Example}

%% ---------- Notation macros ----------
\newcommand{\F}{\mathcal{F}}          % flight set
\newcommand{\Pc}{\mathcal{P}}         % aircraft set
\newcommand{\A}{\mathcal{A}}          % airport set
\newcommand{\D}{\mathcal{D}}          % day set
\newcommand{\C}{\mathcal{C}}          % maintenance check types
\newcommand{\MA}{\mathcal{M}}         % maintenance airport set
\newcommand{\Fl}[2]{F^{\downarrow}_{#1}(#2)}
\newcommand{\Fd}[2]{F^{\uparrow}_{#1}(#2)}
\newcommand{\aI}[1]{a^{\downarrow}_{#1}}
\newcommand{\aD}[1]{a^{\uparrow}_{#1}}
\newcommand{\tI}[1]{t^{\downarrow}_{#1}}
\newcommand{\tD}[1]{t^{\uparrow}_{#1}}
\newcommand{\dt}{\Delta t}

%% highlight corrections
\newcommand{\corrected}[1]{\textcolor{blue}{#1}}
\newcommand{\extended}[1]{\textcolor{teal}{#1}}

%% ---------- Journal metadata (fill before submission) ----------
\journal{European Journal of Operational Research}

\begin{document}

\begin{frontmatter}

\title{Revisiting the Compact Tail Assignment Model: Feasibility and Type-A Maintenance}

\author[lgi]{Viktor Borodin}
\author[lip6]{Arun Kumar Pappala}

\address[lgi]{Taras Shevchenko National University, Kyiv, Ukraine}
\address[lip6]{Singapore University of Technology and Design, Singapore}

\begin{abstract}
This paper revisits the compact mixed-integer linear programming model for
the Tail Assignment Problem (TAP) introduced by Khaled et al.\ (2018), with
scope restricted to flight assignment, aircraft continuity, and type-A
flight-hour maintenance. We examine a simultaneous-departure corner case in
the basic model and formulate explicit pairwise non-overlap constraints to
exclude physically impossible assignments. For maintenance, we distinguish
the interval condition imposed by Constraint~(13) from the opening-state
flight-hour limit required in a finite planning horizon. We describe the
type-A formulation and evaluate its implementations on generated instances,
reporting solver status and independent schedule-validation outcomes. The
generated instances are synthetic and are not derived from the real airline
data set used to generate the test instances
in the original computational study; the results are therefore presented as
implementation experiments, not as a reproduction of that study.
\end{abstract}

\begin{keyword}
Tail assignment \sep Aircraft routing \sep Type-A maintenance \sep
Integer programming
\end{keyword}

\end{frontmatter}

%% ==========================================================================
%%==========================================================================
\section{Introduction}
\label{sec:intro}
%%==========================================================================

The airline planning process is traditionally decomposed into a sequence of
sub-problems: timetable design, fleet assignment, crew rostering, tail
assignment, and recovery planning~\citep{barnhart2003applications,
gopalan1998mathematical}.  The \emph{tail assignment problem} (TAP)—assigning
a specific aircraft (identified by its tail number) to each flight leg—is the
penultimate planning step and is tightly coupled with aircraft maintenance
scheduling.  Poor tail assignment decisions directly inflate operational costs
through unnecessary ferry (repositioning) flights, suboptimal maintenance
routing, and safety-margin violations.

In addition to the basic assignment and continuity requirements, the TAP must
respect aircraft maintenance constraints.  The most common maintenance type is \emph{type-A}, which requires a maintenance check after a maximum flying-time interval.  Type-A maintenance is typically performed at a limited set of
airports, and the TAP must ensure that each aircraft is routed to a maintenance-capable airport before its flying-time limit is exceeded.  The type-A maintenance requirement is a key operational constraint in the TAP and has been studied extensively in the literature~\citep{clarke1997aircraft, barnhart1998flight, gronkvist2005tail}.

\citet{khaled2018compact} proposed a compact 0--1 linear programming
formulation for the TAP without maintenance, featuring only $O(n_f n_p)$ binary variables (where
$n_f$ is the number of flights and $n_p$ the fleet size).  This is a
significant reduction over the classic multi-commodity flow model of
\citet{levin1971scheduling}, which requires $O(n_f^2 n_p)$ variables.
Further, \citet{khaled2018compact} extended the model to include type-A maintenance
requirements, which are defined by a maximum flying-time interval between maintenance checks, that adds other $O(n_f n_p D)$ 
binary variables, where $D$ is the number of calendar days in the planning horizon. 
Their basic-model experiments include instances with about 1\,500 flight legs
and 40 aircraft. Table~6 reports a nonzero residual gap for its largest
no-presolve cell, so these results do not certify every such instance within
one hour. The maintenance experiments in Tables~10--11 are separate and
include three-hour time limits with nonzero residual gaps.

\subsection*{Our contributions}

The present paper makes three contributions that check and correct
those results:

\begin{enumerate}
\item \textbf{Basic feasibility correction (C1--C4 / Khaled constraints (3)--(6)).}
  We show that the strict-time balance definition used in the compact basic
  model admits a simultaneous-departure corner case: one aircraft can be
  assigned to two flights departing at the same airport and same time without
  violating the balance inequalities.  We formalize this counter-example and
  enforce physical executability by adding explicit pairwise non-overlap
  constraints (our C5).
  
  Our non-overlap correction retains the $O(n_fn_p)$ binary-variable space but
can add $O(n_f^2n_p)$ conflict inequalities in the worst case. In practical
timetables, only temporally conflicting pairs are instantiated; the resulting
constraint count is therefore governed by the sparse conflict set
$\mathcal O$ defined in \Cref{sec:basic}. \Cref{subsec:c5_impact} measures
how often the corner case occurs in the benchmark timetables and when it
changes an optimum.


\item \textbf{Cumulative-flight-time modeling.} Constraint~(13) in
  \citet{khaled2018compact} is a valid conditional bound between successive
  checked endpoints, but it is not by itself a complete finite-horizon
  maintenance-hour model: prior exposure and the opening interval to the first
  check need separate protection. We demonstrate an opening gap in the
  implementation's companion row, test a two-row endpoint split on its own
  counterexample, and give exact state formulations for the stated
  day-boundary policy.

\item \textbf{Event-based reformulation.}  The day-indexed type-A model
  above uses calendar-day maintenance decisions. We study two exact encodings
  of an explicitly defined flexible-start, hour-limited extension: one uses
  day-indexed trigger decisions and the other event-indexed decisions with
  actual-start-day selectors. Both allow cross-day deferral and impose one
  check per aircraft per actual start day. They share an event-aware hour
  engine, and their projection equivalence follows by substitution of the
  event aggregate. This is not a fully day-free model, nor an exact
  reformulation of Khaled's complete nightly and calendar-spacing policy.
  A fixed-night restriction is compared separately under otherwise identical
  assumptions. An independent schedule replay checks all loaded incumbents in
  the uniform study; the historical study's replay coverage is reported
  separately in \Cref{sec:computational}.
\end{enumerate}

\subsection*{Paper organisation}

\Cref{sec:literature} reviews related work.  \Cref{sec:basic} restates the
basic TAP model and proves a C1--C4 corner-case correction closed by C1--C5.  \Cref{sec:maintenance}
presents the maintenance-extended model, the cumulative-flight-time row correction, and
the induced interpretation updates for C8--C13 under the corrected basic model.
\Cref{sec:event_reformulation} gives event-aware A-hour constraints and the
uniform flexible encodings, with their projection-equivalence proof.
\Cref{sec:computational} reports the numeric cumulative-flight-time
counterexample, the empirical effect of the non-overlap correction, the
uniform flexible-start comparison, the five-method A-only benchmark with its
replication subsets, and a cross-solver check with Gurobi.
\Cref{sec:conclusion} concludes.

%%==========================================================================
\section{Literature Review}
\label{sec:literature}
%%==========================================================================

\subsection{Mathematical programming formulations}

Three main modelling paradigms underlie TAP research:

\paragraph{Set partitioning / column generation.}
The non-compact formulation enumerates \emph{strings} (feasible flight
sequences per aircraft) as 0--1 variables.  Because the number of strings is
exponential, column generation is required
\citep{lavoie1988new, minoux1984column, barnhart1998flight, sarac2006branch}.
Combining column generation with branch-and-bound (\emph{branch-and-price})
gives exact solutions but at significant algorithmic overhead.  Robust
extensions appear in \citet{borndorfer2010robust} and
\citet{froyland2013recoverable}.

\paragraph{Multi-commodity flow.}
\citet{levin1971scheduling} introduced the compact multi-commodity flow
formulation with $O(n_f^2 n_p)$ binary variables.  This remains the standard
benchmark for compact models and is the primary comparison point in
\citet{khaled2018compact}.

\paragraph{Time-space networks.}
\citet{hane1995fleet} proposed the time-space network for the \emph{fleet
assignment problem}.  Nodes represent airport--time events; arcs encode
ground waits, flights, and overnight stays.  Extensions to aircraft routing
appear in \citet{clarke1996maintenance}.  Preprocessing by node aggregation
and island isolation reduces model size
\citep{hane1995fleet, sherali2006airline}.  Daily and weekly maintenance
routing variants are studied in \citet{liang2011new} and
\citet{liang2012network}; the latter also integrates fleet assignment
decisions, but only approximate solutions are obtained via variable fixing.

\paragraph{Nonlinear / lifted formulations.}
\citet{haouari2012lifted} present a compact formulation with nonlinear
constraints that are linearised via the Reformulation--Linearisation Technique
(RLT).  The largest instances solved involve 344 flights, an order of magnitude
below the roughly 1,500-flight instances reported by
\citet{khaled2018compact}; this is about a fourfold difference in flight
count, not an order of magnitude or a like-for-like comparison of difficulty.

\subsection{MILP alternatives to the Khaled compact model}

Beyond the compact binary assignment model of \citet{khaled2018compact},
several mixed-integer alternatives are relevant for exact TAP modelling and
large-scale implementation:
\begin{itemize}
	\item \textbf{Arc-flow MILP (time-indexed or event-indexed):} decision
	variables are attached to feasible flight-connection arcs instead of flight--aircraft
	assignment pairs, often yielding stronger network-structure constraints
	\citep{levin1971scheduling, hane1995fleet, sherali2006airline}.
	\item \textbf{Route-based master MILP (set partitioning):} aircraft duties are
	represented as columns in a master MILP, typically solved with
	branch-and-price; this is exact but algorithmically heavier
	\citep{barnhart1998flight, sarac2006branch, klabjan2005large}.
	\item \textbf{Decomposed MILP frameworks:} Benders- or
	Dantzig--Wolfe-style decompositions separate assignment/routing from
	maintenance-resource feasibility, improving scalability on larger fleets
	\citep{cordeau2001benders, mercier2005computational, desaulniers1997daily}.
	\item \textbf{Two-stage or rolling-horizon MILP:} first assign flights, then
	repair/optimize maintenance and turnaround feasibility in a second MILP,
	trading global optimality for tractable runtime
	\citep{clarke1996maintenance, liang2011new, liang2012network}.
	\item \textbf{Robust/scenario MILP variants:} delay and disruption uncertainty
	is modelled through multiple scenarios with non-anticipativity-style linking
	constraints, producing schedules that are less brittle operationally
	\citep{borndorfer2010robust, froyland2013recoverable}.
\end{itemize}

In this taxonomy, the Khaled model is best viewed as a compact, directly
solvable baseline that is simple to implement and benchmark, while the
alternatives above trade model size, decomposition complexity, and robustness
for stronger scalability or richer operational realism.

\subsection{Comparison scope and fairness}

For an empirical comparison to be methodologically fair, competing models
should be evaluated under a common instance family, objective definition,
maintenance semantics, and solver-budget protocol.  In TAP, this condition is
often not met across compact MILP, branch-and-price, Benders-decomposed, and
robust/scenario formulations because these approaches typically optimize
different decision scopes or uncertainty assumptions.  Accordingly, this paper
uses \citet{khaled2018compact} as its compact-model reference, not as a
literal reproduction of the complete original maintenance policy. The main
instance-matched study compares uniform flexible encodings; historical
C13 controls and fixed-night timing are labelled separately. Other MILP
families are complementary paradigms rather than one-to-one benchmarks.

\subsection{Maintenance constraints}

Maintenance scheduling within aircraft routing has been studied since the
early work of \citet{clarke1997aircraft}. The column-generation framework
of \citet{barnhart1998flight} can embed type-A constraints through constrained
shortest-path pricing.  \citet{cordeau2001benders} apply Benders decomposition
to simultaneous routing and crew scheduling with non-linear maintenance
resources.  \citet{lacasse2010aircraft} consider maintenance under different
business-process models. We focus on the type-A flight-hour case and its exact
event-based reformulation.

\subsection{Heuristic approaches}

Constructive heuristics for TAP are comparatively under-studied relative to
exact methods.  \citet{gronkvist2005tail} provides a comprehensive treatment
including local-search improvements. Heuristic and metaheuristic solution
methods are outside the scope of the present exact-formulation study and are
not evaluated here.

\subsection{Position of the present paper}

The compact formulation of \citet{khaled2018compact} represents the
reference compact baseline for exact, cyclic-constraint-free TAP in this
study.  The present work builds directly upon it, providing: (i) a
diagnosis of the C1--C13 formulation in \citet{khaled2018compact} and an exact state-based replacement for the legacy cumulative-flight-time row, (ii) a
uniform flexible-start A-check encodings with an exact hour engine, and
(iii) an instance-matched benchmark that separates encoding choices from
maintenance-policy restrictions. Computational superiority is an empirical
question, not a property assumed by the formulation.

%%==========================================================================
\section{The Basic Tail Assignment Model}
\label{sec:basic}
%%==========================================================================

We restate the compact model of \citet{khaled2018compact} using the notation
of \Cref{tab:notation}.  This section covers the model without maintenance
constraints; they are introduced in \Cref{sec:maintenance}.

\begin{table}[t]
\centering\small
\caption{Notation summary.}
\label{tab:notation}
\begin{tabular}{ll}
\toprule
Symbol & Meaning \\
\midrule
$\F = \{1,\ldots,n_f\}$     & Set of flight legs \\
$\Pc = \{1,\ldots,n_p\}$    & Set of aircraft (tail numbers) \\
$\A = \{1,\ldots,n_k\}$     & Set of airports \\
$\aD{i}$, $\aI{i}$          & Departure / arrival airport of flight $i$ \\
$\tD{i}$, $\tI{i}$          & Departure / arrival time of flight $i$ (minutes) \\
$d_i$, $\bar d_i$           & Departure / arrival day derived from flight times \\
$\dt$                        & Minimum ground turn time (minutes) \\
$c_{ij}$                     & Assignment cost of flight $i$ to aircraft $j$ \\
$a^0_j$                      & Initial position airport of aircraft $j$ \\
$\Fl{k}{\theta}$             & $\{i \in \F : \aI{i} = k,\; \tI{i} \le \theta - \dt\}$ \\
$\Fd{k}{\theta}$             & $\{i \in \F : \aD{i} = k,\; \tD{i} < \theta\}$ \\
\bottomrule
\end{tabular}
\end{table}

\subsection{Decision variables and objective}

Let us define the decision variables and constraints for the basic tail assignment model.
\paragraph{Binary assignment.}
The decision variable $x_{ij}$ indicates whether flight $i$ is assigned to aircraft $j$.
\begin{equation}\label{eq:c0}
  x_{ij} = \begin{cases} 1 & \text{if flight } i \text{ is operated by aircraft } j, \\ 0 & \text{otherwise,} \end{cases}
  \qquad i \in \F,\; j \in \Pc.
\end{equation}

The objective minimises total assignment cost:
\begin{equation}\label{eq:obj}
  \min \sum_{i \in \F} \sum_{j \in \Pc} c_{ij} x_{ij}.
\end{equation}

\subsection{Constraints}

\paragraph{C2 — Flight coverage.}
Each flight must be assigned to exactly one aircraft:
\begin{equation}\label{eq:c1}
  \sum_{j \in \Pc} x_{ij} = 1, \qquad \forall i \in \F.
\end{equation}

\paragraph{C3,C4 — Equipment continuity (home/non-home airports).}
For any aircraft $j$, any non-home airport $k \ne a^0_j$, and any flight
$i$ departing from $k$, aircraft $j$ can be assigned to flight $i$ only if
it has already landed at $k$ early enough,
but at the home airport ($k = a^0_j$), one free unit is available at time zero
\begin{equation}\label{eq:c2}
  \sum_{i' \in \Fl{k}{\tD{i}}} x_{i'j}
  - \sum_{i' \in \Fd{k}{\tD{i}}} x_{i'j}
  \;\geq\; x_{ij},
  \qquad \forall j \in \Pc,\; k \in \A \setminus \{a^0_j\},\;
  i \in \F : \aD{i}=k.
\end{equation}

\begin{equation}\label{eq:c3}
  \sum_{i' \in \Fl{k}{\tD{i}}} x_{i'j}
  - \sum_{i' \in \Fd{k}{\tD{i}}} x_{i'j}
  \;\geq\; x_{ij}-1,
  \qquad \forall j \in \Pc,\; k = a^0_j,\;
  i \in \F : \aD{i}=k.
\end{equation}

\subsection{Validity}

Constraints~\eqref{eq:c2}--\eqref{eq:c3} are intended to encode both equipment
continuity (C2 of \citet{khaled2018compact}) and turn-time feasibility
(constraint~(6) of that paper).
The set $\Fl{k}{\tD{i}}$ counts only arrivals at least $\dt$ minutes before
the departure of flight $i$, so a prior arrival can supply flight $i$ only
after the required turn time. This mechanism is valid when departures are
strictly ordered, but it does not prevent two equal-time departures from using
the same available aircraft, as shown below.

\begin{proposition}[Claimed in \citealt{khaled2018compact}]
  Any feasible solution to \eqref{eq:c0}--\eqref{eq:c3}
  corresponds to a
  feasible flight plan satisfying the original continuity and turn-time
  conditions.
\end{proposition}

The proposition is not valid as stated when distinct flights may have equal
departure times. The following counterexample identifies the missing case.

\subsection{Counter-example with simultaneous departures}

The balance constraints~\eqref{eq:c2}--\eqref{eq:c3} use
$\Fd{k}{\theta}=\{i\in\F:\aD{i}=k,\;\tD{i}<\theta\}$ with a strict
inequality in departure time. If two flights leave the same airport at the
same time, neither flight appears in the other's departure-balance term.
This can make an assignment feasible while being physically impossible as a
flight plan.

\begin{example}[Feasible to \eqref{eq:c0}--\eqref{eq:c3}, infeasible in practice]
  \label{ex:simultaneous_departure_counterexample}
  Consider one aircraft $j$ with initial airport $a^0_j=\text{B}$ and
  turn time $\dt=30$ minutes. Let the three flights be:
  \begin{align*}
    &i_0: \text{B}\to\text{A},\quad \tD{i_0}=500,\;\tI{i_0}=550,\\
    &i_1: \text{A}\to\text{B},\quad \tD{i_1}=600,\;\tI{i_1}=700,\\
    &i_2: \text{A}\to\text{C},\quad \tD{i_2}=600,\;\tI{i_2}=700.
  \end{align*}
  Set $x_{i_0j}=x_{i_1j}=x_{i_2j}=1$.

  Coverage~\eqref{eq:c1} is satisfied (single aircraft, all flights assigned).
  Flight $i_0$ satisfies the home-airport constraint~\eqref{eq:c3}, because
  no flight has arrived at or departed from B before time 500. For $i_1$ at
  the non-home airport A:
  \begin{equation*}
    \Fl{\text{A}}{\tD{i_1}}=\{i_0\},\qquad
    \Fd{\text{A}}{\tD{i_1}}=\varnothing,
  \end{equation*}
  because $i_0$ departs from B and $\tD{i_2}=600\not<600$. Hence
  constraint~\eqref{eq:c2} gives
  \begin{equation*}
    \sum_{i'\in\Fl{\text{A}}{\tD{i_1}}}x_{i'j}
    -\sum_{i'\in\Fd{\text{A}}{\tD{i_1}}}x_{i'j}
    =1-0=1
    \ge x_{i_1j}=1.
  \end{equation*}
  Symmetrically for $i_2$, $\tD{i_1}=600\not<600$, so the same calculation gives
  $1\ge x_{i_2j}=1$. Therefore~\eqref{eq:c0}--\eqref{eq:c3} are feasible.

  However, the resulting plan is physically infeasible: one aircraft cannot
  operate both $i_1$ and $i_2$ simultaneously at time 600.
\end{example}

\paragraph{Implication.}
In implementations, this corner case is eliminated by adding pairwise
non-overlap constraints for time-overlapping flights on the same aircraft,
$x_{i_1j}+x_{i_2j}\le 1$.

\paragraph{Corrected C1--C5 feasibility criterion.}
Define $\mathcal{O}$ as the set of flight pairs that cannot be operated by one
aircraft because their operating intervals, including the turn-time buffer,
overlap:
\paragraph{Conflict-set definition.}
\begin{equation}
  \mathcal{O}=\left\{\{i_1,i_2\}\subseteq\F:
  [\tD{i_1},\tI{i_1}+\dt)\cap
  [\tD{i_2},\tI{i_2}+\dt)\neq\varnothing\right\}.
  \label{eq:overlap_set}
\end{equation}
For every pair in this conflict set, impose
\paragraph{C5 — Pairwise non-overlap.}
\begin{equation}\label{eq:c4_overlap}
  x_{i_1 j} + x_{i_2 j} \le 1,
  \qquad \forall \{i_1,i_2\}\in\mathcal{O},\; \forall j\in\Pc.
\end{equation}

The corrected core model is therefore
\eqref{eq:c0}--\eqref{eq:c4_overlap}.
This preserves the original $n_fn_p$ binary-variable space while restoring
physical executability. It does not, in general, preserve the linear
constraint count: the overlap family contains $|\mathcal{O}|n_p$ inequalities.

\subsection{Model size}

\begin{proposition}[Analogue of \citealt{khaled2018compact}, Proposition~2]
  The original model \eqref{eq:obj}--\eqref{eq:c3} with binary $x$ contains $n_f n_p$ binary
  decision variables and $O(n_fn_p)$ constraints, while the corrected model has $n_fn_p$ binary variables and $O(n_fn_p+|\mathcal{O}|n_p)$ constraints.
\end{proposition}

The corrected model retains $n_fn_p$ binary variables and has
$O(n_fn_p+|\mathcal{O}|n_p)$ constraints. In the worst case
$|\mathcal{O}|=O(n_f^2)$, although timetable sparsity normally makes the
conflict set much smaller. The original variable count compares favourably
with Levin's $O(n_f^2n_p)$-variable model; \citet{khaled2018compact} also
report that the time-space network formulation uses approximately
$1.25$--$2$ times as many variables and constraints on their test instances.

\subsection{Feasibility of the corrected basic model}

\begin{proposition}
  Any binary solution to \eqref{eq:c1}--\eqref{eq:c3} and
  \eqref{eq:c4_overlap} defines a physically executable flight sequence for
  each aircraft, under the stated deterministic turn-time assumptions.
\end{proposition}

\begin{proof}
  Fix an aircraft $j$ and order its assigned flights by departure time. At
  each assigned departure from airport $k$, constraints~\eqref{eq:c2}--
  \eqref{eq:c3} require at least one unit of aircraft availability at $k$,
  computed from the initial position and arrivals that meet the turn-time
  requirement, less departures from $k$ at strictly earlier times. Thus each
  departure can be reached from the initial position or from a previously
  operated flight. Constraint~\eqref{eq:c4_overlap} excludes any pair of
  assigned flights whose operating intervals, including turn time, overlap.
  Applying these facts in chronological order yields a continuous executable
  sequence for aircraft $j$. Since $j$ was arbitrary, the result holds for
  every aircraft.
\end{proof}







%%==========================================================================
\section{Maintenance Extension and Cumulative-Hour Coverage}
\label{sec:maintenance}
%%==========================================================================

\subsection{Maintenance problem definition}

\paragraph{Primary-source scope.}
Section~5.1 and Table~9 of \citet{khaled2018compact} (p.~555) specify
night maintenance at the arrival airport of the last flight of the day,
capacity measured as aircraft handled per airport-day, and both calendar
spacing and cumulative-hour requirements for type-A checks. Their
equation~(13) sums over $F_{d+1,d'}$, i.e. the flights departing on days
$d+1,\ldots,d'$; with the convention in this section, this is
$F_{d,d'}=(d,d']$.
These definitions do not describe arbitrary flexible daytime starts.

The main experiment here is an explicitly different, hour-limited A-check
extension: starts may be deferred across days, checks have a stated duration,
and capacity measures simultaneous bay occupancy. It omits the original
calendar-spacing requirement rather than assigning that requirement only to
C/D checks. All main encodings use these same assumptions. Fixed-night
starts are a separate restricted comparison, not a relaxation of flexible
timing. The following day-indexed rows provide the historical and
day-boundary reference for the cumulative-hour analysis; they are not the
complete flexible-start model.

Additional notation is collected in \Cref{tab:maint_notation}.

\begin{table}[t]
\centering\small
\caption{Additional notation for the maintenance model.}
\label{tab:maint_notation}
\begin{tabularx}{\linewidth}{lX}
\toprule
Symbol & Meaning \\
\midrule
$\D = \{1,\ldots,n_d\}$      & Planning days \\
$d_i$                         & Departure day of flight $i$ \\
$\bar d_i$                    & Arrival day of flight $i$ \\
$\MA \subseteq \A$            & Maintenance-capable airports \\
$\text{mcap}_{m,d}$           & Maintenance capacity of airport $m$ on day $d$ \\
$mc_{ijd}$                    & Cost of night maintenance at arrival airport of
                                flight $i$, aircraft $j$, day $d$ \\
$F_d^{\downarrow}(m)$         & Flights arriving at airport $m$ on day $d$ \\
$F(d,i)$                      & $\{i' : i'\ne i,\; d_{i'}=d,\; \tD{i'} \geq \tI{i}\}$ \\
$T^A_{\max}$                  & Max cumulative flight minutes before type-A check \\
$d_{\max}$                    & Max calendar-day interval between checks \\
$M_{dd'}$                     & Big-$M$ for day pair $(d,d')$ \\
$F_{d,d'}$                    & Flights with departure day in $(d,d']$ \\
\bottomrule
\end{tabularx}
\end{table}

\subsection{Additional decision variables}
\paragraph{C6, C7 — Boolean variables for maintenance.}
\begin{align*}
  z_{ijd} &= \begin{cases}1 & \text{night maintenance at arrival airport of } i,\;
               \text{aircraft } j,\;\text{day } d, \\ 0 & \text{else;}\end{cases}\\
  y_{jd}  &= \begin{cases}1 & \text{maintenance of aircraft } j \text{ on day } d,
               \\ 0 & \text{else,}\end{cases}
\end{align*}

The variable $z_{ijd}$ is instantiated only when flight $i$ arrives at a
maintenance-capable airport on day $d$. Thus its admissible domain is
$i\in\bigcup_{m\in\MA}F_d^{\downarrow}(m)$.

\paragraph{Overnight executability assumption.}
Let $\tau_A>0$ be the check duration and let $[L_{md},U_{md}]$ be the
available service window at maintenance airport $m$ on the night following
day $d$, measured in the same absolute time units as the flight timestamps.
Interpret $\mathrm{mcap}_{m,d}$ as the number of independent, staffed check
slots reserved for that night, not merely as a daily throughput quota. For
every selected trigger $z_{ijd}=1$, assume there exists a check start time
$s_{ijd}$ such that
\begin{align}
  s_{ijd}&\geq \tI{i}, &
  L_{md}&\leq s_{ijd}, &
  s_{ijd}+\tau_A&\leq U_{md},
  \label{eq:c_night_window}
\end{align}
and, for every distinct flight $i'\ne i$ assigned to aircraft $j$ that departs
at or after the trigger arrival, $\tD{i'}\geq s_{ijd}+\tau_A+\dt$. C8
excludes same-day departures at or after the trigger arrival; this last
condition also covers the first assigned
departure on a subsequent day. If these conditions are not guaranteed by the
input timetable and operating policy, the constraints listed here alone do
not establish check completion before the aircraft's next departure.

\subsection{Extended model}

The objective becomes:
\begin{equation}\label{eq:obj_ext}
  \min\; \sum_{i \in \F}\sum_{j \in \Pc} c_{ij} x_{ij}
  + \sum_{d \in \D}\sum_{m \in \MA}\sum_{i \in F_d^{\downarrow}(m)}
    \sum_{j \in \Pc}
    mc_{ijd}\, z_{ijd}.
\end{equation}

Constraints~\eqref{eq:c0}--\eqref{eq:c4_overlap} are retained.  Additional
constraints:

\paragraph{C8 — Maintenance blocks same-day subsequent flights.}
\begin{equation}\label{eq:c8}
  z_{ijd} + x_{i'j} \le 1,
  \quad \forall d\in\D,\;m\in\MA,\;i\in F_d^{\downarrow}(m),\;
  i'\in F(d,i),\;j\in\Pc.
\end{equation}

\paragraph{C9 — Maintenance requires assignment.}
\begin{equation}\label{eq:c9}
  x_{ij} \ge z_{ijd},
  \quad \forall d\in\D,\;m\in\MA,\;i\in F_d^{\downarrow}(m),\;j\in\Pc.
\end{equation}

\paragraph{C10 — Capacity.}
\begin{equation}\label{eq:c10}
  \sum_{i \in F_d^{\downarrow}(m)} \sum_{j \in \Pc} z_{ijd}
  \le \text{mcap}_{m,d},
  \qquad \forall d \in \D,\; m \in \MA.
\end{equation}

\paragraph{C11 — Aggregation ($z \to y$).}
\begin{equation}\label{eq:c11}
  \sum_{m\in\MA}\sum_{i \in F_d^{\downarrow}(m)} z_{ijd} = y_{jd},
  \qquad \forall d \in \D,\; j \in \Pc.
\end{equation}

\paragraph{C12 — Maximum-spacing constraint.}
\begin{equation}\label{eq:c12}
  \sum_{r \in [d, d']} y_{jr} \ge 1,
  \qquad \forall d, d' \in \D \text{ s.t.\ } d' - d = d_{\max} - 1,\;
  \forall j \in \Pc.
\end{equation}

Khaled's type-A model includes C12. Our hour-only experiments deliberately
leave it inactive and must therefore not be described as a reproduction of
all original type-A requirements. An experiment imposing C12 would also need
an explicitly supplied initial calendar age and common opening-window policy.

The original paper's cumulative-flight-time constraint is C13:
\paragraph{C13 — Paper-based cumulative-flight-time constraint.}
\begin{equation}\label{eq:c13_paper}
  \sum_{i \in F_{d,d'}} x_{ij}\,\tilde{t}_i
  \;\le\;
  T^A_{\max}
  + M_{dd'}\Bigl(2 - y_{jd} - y_{jd'}\Bigr)
  + M_{dd'}\!\sum_{r \in [d+1,d'-1]} y_{jr},
\end{equation}
for all $j \in \Pc$, all pairs $(d, d')$ with $2 \le d'-d \le d_{\max}-1$,
$d < n_d$, where $\tilde{t}_i$ is the flight duration of leg $i$.


\subsection{Impact of corrected C1--C5 feasibility on C8--C13}
\label{subsec:c2c4_impact_maintenance}

The correction introduced in \Cref{sec:basic} (non-overlap
constraint~\eqref{eq:c4_overlap}) changes the admissible assignment set for
$x$, and therefore the interpretation of all maintenance constraints layered on
top of $x$.

\paragraph{C8.}
No algebraic change is required in~\eqref{eq:c8}; however, with
\eqref{eq:c4_overlap} active, the blocking relation now operates on physically
realizable trajectories only. This removes spurious cases where maintenance
placement was evaluated on infeasible simultaneous departures.

\paragraph{C9--C12.}
Constraints~\eqref{eq:c9}--\eqref{eq:c12} remain structurally unchanged. Their
feasible region cannot expand because $x$ is restricted by C5. It is strictly
smaller only when the original rows admitted an assignment that uses at least
one overlapping flight pair on the same aircraft.

\paragraph{C13 and cumulative-hour rows.}
Hour-accumulation constraints are evaluated over a physically executable
$x$-space once~\eqref{eq:c4_overlap} is imposed. This correction is necessary,
but it does not by itself establish that a particular cumulative-hour algebraic form
correctly propagates the maintenance counter; that property is analysed below.

In summary, the C1--C5 correction primarily changes the domain over which the
maintenance constraints operate. It is independent of the separate algebraic
correction required for cumulative-hour propagation.

\subsection{Issues of day-indexed maintenance constraints C8--C13}

Usage of day-indexed maintenance variables $y_{jd}$ and $z_{ijd}$ is
operationally intuitive, gives clear modeling semantics, but it has the following drawbacks:
\begin{itemize}
\item The day-indexed formulation requires a large number of big-$M$ constants
  $M_{dd'}$ to cover all day pairs $(d,d')$ with $2\le d'-d\le d_{\max}-1$.
  The number of such pairs is polynomial, at most
  $O(n_d\min\{n_d,d_{\max}\})$. Each $M_{dd'}$ must bound the possible
  excess only for its own flight window $F_{d,d'}$; its required magnitude
  depends on that window and on the available assignment bounds, not on all
  flight hours over the entire horizon.

  \item The C13 rows do not propagate a flight-hour state. Each row bounds
  one day window only under a particular endpoint pattern, so exposure that
  falls outside every active window is unbounded; the five resulting gaps are
  analysed in \Cref{subsec:c13_opening_coverage}. An exact day-indexed state
  recursion does exist and is given in \Cref{subsec:c11_state}.

  \item A calendar-day limit is not a rolling 24-hour limit. Checks on either
  side of midnight may be less than 24 hours apart; two starts on the same
  calendar day are disallowed in the uniform study.

  \item Ordered flight states offer an alternative representation without
  successor variables. Their total size and runtime depend on event conflicts
  and do not imply computational superiority over day-indexed decisions.

  \item C13 is conditional on successive checked endpoints and is indexed
  only for $2\le d'-d\le d_{\max}-1$. It contains no pre-horizon exposure and
  no row for the opening period, for consecutive-day checks, or for checks
  $d_{\max}$ days apart, and it relaxes every closing window.
\end{itemize}


\subsection{Coverage gaps of the published cumulative-flight-time row}
\label{subsec:c13_opening_coverage}

Throughout this subsection, checks follow the flights of their labelled day,
as in C8. Write $v_{jd}$ for the flying time of aircraft $j$ on day $d$
(formalised in~\eqref{eq:daily_flight_time}),
$L_{jdd'}=\sum_{r=d+1}^{d'}v_{jr}$ for the left-hand side of C13, and
$I_{jdd'}=\sum_{r=d+1}^{d'-1}y_{jr}$ for its interior check count.

\paragraph{Physical rule.}
The accumulated exposure of aircraft $j$ immediately before the possible
check of night $d$ is
\begin{equation}\label{eq:physical_exposure}
  H_{jd}=\begin{cases}
    \Phi^A_j+\sum_{r=1}^{d}v_{jr}, & \text{if } y_{jr}=0 \text{ for all } r<d,\\[2pt]
    L_{jsd}, & \text{otherwise, where } s=\max\{r<d:\ y_{jr}=1\}.
  \end{cases}
\end{equation}
The type-A hour limit requires $H_{jd}\le T^A_{\max}$ for every aircraft and
every day $d\in\D$. Flights of a check day count toward the exposure before
that check.

\begin{lemma}[Validity of the published and split rows]
  \label{lem:c13_validity}
  Let $y$ be binary, $v\ge0$, and let every $M_{dd'}$ be at least the largest
  possible excess of $L_{jdd'}$ over $T^A_{\max}$. Every schedule satisfying
  the physical rule satisfies~\eqref{eq:c13_paper} and the split
  rows~\eqref{eq:c11_split_a}--\eqref{eq:c11_split_b} below.
\end{lemma}

\begin{proof}
  A row whose right-hand side contains a positive multiple of $M_{dd'}$ is
  satisfied by the choice of $M_{dd'}$. Each of the three rows has right-hand
  side $T^A_{\max}$ only if $I_{jdd'}=0$. Then no check occurs on days
  $d+1,\ldots,d'-1$, so the last check before day $d'$, if any, occurs on or
  before day $d$. Since $v\ge0$, \eqref{eq:physical_exposure} gives
  $L_{jdd'}\le H_{jd'}\le T^A_{\max}$.
\end{proof}

Thus C13 never cuts off a physically valid schedule; its defect is
incomplete coverage. The physical rule must be bounded on day $d$ by some
active row, but C13 is active only for a window $(d,d']$ that satisfies all
of the following: both endpoints are checks, there is no interior check,
$2\le d'-d\le d_{\max}-1$, and $d<n_d$. Every exposure not of this form is
unbounded by C13.

\begin{proposition}[Coverage gaps of Constraint~(13)]
  \label{prop:c13_gaps}
  Even together with calendar spacing~\eqref{eq:c12}, the published rows do
  not bound the following exposures:
  \begin{enumerate}
    \item[G1] pre-horizon exposure $\Phi^A_j$, which appears in no row;
    \item[G2] the opening period before the first check, because every row
      with $y_{jd}=0$ at its start is relaxed and no row starts before day~1;
    \item[G3] consecutive check days, $d'-d=1$, which are outside the index
      range (the paper's prose range $1\le d'-d$ would cover this case, but
      the displayed equation excludes it);
    \item[G4] successive checks exactly $d_{\max}$ days apart, which
      \eqref{eq:c12} permits but the index range $d'-d\le d_{\max}-1$
      excludes;
    \item[G5] the closing period after the last check, because its end day
      is not a check day and the row is relaxed.
  \end{enumerate}
\end{proposition}

\begin{proof}
  \Cref{tab:c13_gap_witnesses} gives, for each gap, daily data with
  $T^A_{\max}=600$ and $n_d=6$ that satisfy~\eqref{eq:c12} and every row
  of~\eqref{eq:c13_paper}, yet violate the physical rule. Each check day
  carries positive flying time, so every check can follow an arrival, for
  example on loops from a single maintenance base. Acceptance holds for every
  valid choice of $M_{dd'}$, because a relaxed row needs only $M_{dd'}$ at
  least its own excess. The entries were verified by evaluating all indexed
  rows and~\eqref{eq:physical_exposure}.
\end{proof}

\begin{table}[t]
\centering\scriptsize
\caption{Witnesses for the coverage gaps of Constraint~(13)
  ($T^A_{\max}=600$, $n_d=6$). ``Accepts'' means that every indexed row is
  satisfied; ``rejects'' means that some row is violated. G0 is the
  check-to-check case on which the endpoint split fails. Exact means the
  state recursion of \Cref{subsec:c11_state} or, equivalently, the
  interval-prefix rows of \Cref{subsec:c13_prefix_equivalence}.}
\label{tab:c13_gap_witnesses}
\resizebox{\linewidth}{!}{%
\begin{tabular}{llrllrllll}
\toprule
Gap & $d_{\max}$ & $\Phi^A_j$ & $(v_{j1},\ldots,v_{j6})$ & $(y_{j1},\ldots,y_{j6})$ &
  Overrun & C13 & Split & C13 + split & Exact \\
\midrule
G0 & 4 & 0   & $(10,600,0,600,0,0)$   & $(1,0,0,1,0,0)$ & $H_{j4}=1200$ & rejects & accepts & rejects & rejects \\
G1 & 3 & 590 & $(0,20,0,10,0,0)$      & $(0,1,0,1,0,0)$ & $H_{j2}=610$  & accepts & accepts & accepts & rejects \\
G2 & 3 & 0   & $(400,400,10,0,0,10)$  & $(0,0,1,0,0,1)$ & $H_{j2}=800$  & accepts & accepts & accepts & rejects \\
G3 & 3 & 0   & $(10,700,0,10,0,0)$    & $(1,1,0,1,0,0)$ & $H_{j2}=700$  & accepts & accepts & accepts & rejects \\
G4 & 3 & 0   & $(10,400,0,410,0,0)$   & $(1,0,0,1,0,0)$ & $H_{j4}=810$  & accepts & accepts & accepts & rejects \\
G5 & 3 & 0   & $(10,0,0,10,400,400)$  & $(1,0,0,1,0,0)$ & $H_{j6}=800$  & accepts & rejects & rejects & rejects \\
\bottomrule
\end{tabular}}
\end{table}

\paragraph{Historical opening companion.}
Gaps G1--G2 suggest adding an opening row. The companion used with C13 in
the historical implementation reads
\begin{equation}\label{eq:c13b_legacy}
  \Phi^A_j+\sum_{r=1}^{t}v_{jr}
  \le T^A_{\max}+M^0\sum_{r=1}^{t}y_{jr},
  \qquad t=2,\ldots,n_d,
\end{equation}
with $M^0$ a valid bound on the opening excess.

\begin{lemma}[The historical opening companion does not close G1--G2]
  \label{lem:c13_opening_gap}
  Constraints~\eqref{eq:c12}, \eqref{eq:c13_paper} and~\eqref{eq:c13b_legacy}
  together admit schedules that violate the physical rule before the first
  check.
\end{lemma}

\begin{proof}
  The companion has two defects. First, its sum includes the end-day indicator
  $y_{jt}$, although the flights of day $t$ precede that night's check. Take
  $d_{\max}=4$, $\Phi^A_j=0$, $v_j=(0,590,0,20,0,0)$ and
  $y_j=(0,0,0,1,0,0)$. Spacing holds, every C13 row has an unchecked endpoint
  and is relaxed, rows $t=2,3$ read $590\le600$, and rows $t\ge4$ are relaxed
  by $y_{j4}$. Nevertheless, $H_{j4}=610>600$. Second, no row exists for
  $t=1$. Take $\Phi^A_j=590$, $v_j=(20,0,0,10,0,0)$ and
  $y_j=(1,0,0,1,0,0)$. Every companion row is relaxed by $y_{j1}$, the only
  active C13 row is $(1,4)$ with $L_{j14}=10$, and $H_{j1}=610>600$.
\end{proof}

Both defects disappear if the end-day indicator is excluded, i.e.\ the sum
on the right runs to $t-1$, and the row is imposed from $t=1$. This is the
opening row~\eqref{eq:c13_prefix_opening}. Together with the interval rows,
it closes G1 and G2 only; G3--G5 are interval gaps and remain under the
published index range.

\subsection{Legacy cumulative-flight-time row: endpoint-split strengthening and limitation}
\label{subsec:c11_corrected}

An endpoint-split alternative considered for replacing
\eqref{eq:c13_paper} consists of two constraints, each relaxed by one endpoint:

\begin{equation}\label{eq:c11_split_a}
  \sum_{i \in F_{d,d'}} x_{ij}\,\tilde{t}_i
  \;\le\;
  T^A_{\max}
  + M_{dd'}\, y_{jd}
  + M_{dd'}\!\sum_{r \in [d+1,d'-1]} y_{jr},
\end{equation}
\begin{equation}\label{eq:c11_split_b}
  \sum_{i \in F_{d,d'}} x_{ij}\,\tilde{t}_i
  \;\le\;
  T^A_{\max}
  + M_{dd'}\, y_{jd'}
  + M_{dd'}\!\sum_{r \in [d+1,d'-1]} y_{jr},
\end{equation}
for all $j \in \Pc$, $(d,d') \in \D^2$ with $2 \le d'-d \le d_{\max}-1$,
$d < n_d$.

\begin{lemma}[The endpoint split is not sufficient]
  \label{lem:c11_split_limit}
  With $F_{d,d'}$ defined as the flights departing in $(d,d']$,
  constraints~\eqref{eq:c11_split_a}--\eqref{eq:c11_split_b} do not necessarily
  bound all flying time between two consecutive maintenance nights.
\end{lemma}

\begin{proof}
  Let $T^A_{\max}=10$, $M_{dd'}=20$, $d_{\max}\geq4$, and consider days
  $1,\ldots,4$ with
  $y_{j1}=y_{j4}=1$ and $y_{j2}=y_{j3}=0$. Assign 10 units of flying time on
  day 2 and 10 units on day 4. For $(d,d')=(1,3)$, the accumulated time is 10
  and the second split inequality is active. For $(2,4)$, the accumulated time
  is 10 and the first split inequality is active. For $(1,4)$, the accumulated
  time is 20, but both inequalities have right-hand side
  $T^A_{\max}+M_{14}=30$ because both endpoints contain checks. Thus every
  split inequality is satisfied although 20 units are flown between the
  consecutive checks, exceeding the limit 10.
\end{proof}

The original and split formulations are therefore not ordered by strength:
the original row is active when both endpoints contain checks, whereas the
split rows are active when at least one corresponding endpoint indicator is
zero. Neither pattern alone gives an exact cumulative counter.

\paragraph{What the split repairs and what it does not.}
For binary $y$ and a common valid $M_{dd'}$, the published row and the two
split rows for the same pair are jointly equivalent to the strict interval row
\begin{equation}\label{eq:c13_strict_pair}
  L_{jdd'}\le T^A_{\max}+M_{dd'}I_{jdd'}.
\end{equation}
Indeed, if $I_{jdd'}\ge1$ all four rows are relaxed. If $I_{jdd'}=0$, the
strict row reads $L_{jdd'}\le T^A_{\max}$; this is the published row when
$y_{jd}=y_{jd'}=1$, the second split row when $y_{jd'}=0$, and the first
split row when $y_{jd}=0$. Conversely, each of the three rows is implied
by~\eqref{eq:c13_strict_pair}. The split alone detects the G5 witness within
the indexed pair range, but does not close the trailing period in general:
it misses a one-day tail, for which no row is indexed, and loses G0, as shown
in \Cref{lem:c11_split_limit}. Even adding the split to C13 yields
strict interval rows only on $2\le d'-d\le d_{\max}-1$. G1 and G2 remain
because neither family contains $\Phi^A_j$ or a row starting at the horizon,
and G3 and G4 remain because their windows are outside the index range
(\Cref{tab:c13_gap_witnesses}). Rearranging endpoint indicators cannot
remove G1--G4: closing G3--G4 requires interval rows for every pair
$d<d'$, and closing G1--G2 requires an opening state that includes
$\Phi^A_j$.

\subsection{Exact cumulative-state formulation for cumulative flight hours}
\label{subsec:c11_state}

Define the flying time assigned to aircraft $j$ on day $d$ by
\paragraph{C13a Daily flight time}
\begin{equation}
  v_{jd}=\sum_{i\in\F:\,d_i=d}\tilde{t}_i x_{ij}.
  \label{eq:daily_flight_time}
\end{equation}
Let $h_{jd}\geq0$ be the accumulated type-A flying time immediately after the
flights of day $d$ and before maintenance that night. Set
$h_{j0}=\Phi^A_j$ and $y_{j0}=0$. For consecutive days, impose
\paragraph{C13b-C13f Accumulated flight time constraints}
\begin{align}
  h_{jd}
  &\geq h_{j,d-1}+v_{jd}-M^H y_{j,d-1},
  \label{eq:c11_state_lower}\\
  h_{jd}
  &\leq h_{j,d-1}+v_{jd}+M^H y_{j,d-1},
  \label{eq:c11_state_upper}\\
  h_{jd}
  &\geq v_{jd},
  \label{eq:c11_state_reset_lower}\\
  h_{jd}
  &\leq v_{jd}+M^H(1-y_{j,d-1}),
  \label{eq:c11_state_reset_upper}\\
  0\leq h_{jd}&\leq T^A_{\max},
  \qquad \forall j\in\Pc,\ d\in\D,
  \label{eq:c11_state_bound}
\end{align}
where $M^H\geq T^A_{\max}$ is valid because
$h_{j,d-1}\leq T^A_{\max}$. If $y_{j,d-1}=0$, the first two inequalities give
$h_{jd}=h_{j,d-1}+v_{jd}$. If $y_{j,d-1}=1$, the reset inequalities give
$h_{jd}=v_{jd}$. The upper bound then enforces the flight-hour threshold on
every day, including the flights operated immediately before a night check.

The state $h_{jd}$ is indexed at day boundaries, whereas the ordered
Event-based state in \Cref{sec:event_reformulation} is indexed after individual
flights and resets at a selected event. These are distinct time resolutions;
no general projection-equivalence theorem between the formulations is claimed
here. Section~\ref{subsec:event_measurements} compares them under the stated
A-check timing policies and independently audits optimizer-selected schedules.

\subsection{Equivalent implementation strengthening}
\label{subsec:corrected_legacy_strengthening}

The exact-state formulation admits a smaller, tighter implementation without
changing its integer feasible region. The strengthened implementation is
type-A-only, matching the scope of this paper. Its state variable receives the
threshold directly as a variable bound, $0\le h_{jd}\le T^A_{\max}$. The explicit copy of this upper bound
in~\eqref{eq:c11_state_bound} is then unnecessary in the generated matrix.

Define the aircraft-day upper bound
\begin{equation}
  \bar v_{jd}=\sum_{i\in\F:\,d_i=d,\,(i,j)\in\mathcal X}\tilde t_i
  \label{eq:strengthened_daily_upper}
\end{equation}
and the reachable state bounds
\begin{equation}
  U_{j0}=\min\{T^A_{\max},\Phi^A_j\},\qquad
  U_{jd}=\min\{T^A_{\max},U_{j,d-1}+\bar v_{jd}\}.
  \label{eq:strengthened_state_upper}
\end{equation}
Because $h_{jd}\le U_{jd}$ is valid by induction, the reset linearization for
day $d$ may use $U_{j,d-1}$ instead of the horizon-wide
$T^A_{\max}$. This is particularly tighter near the beginning of the
planning horizon.

Finally, the strengthened implementation generates maximal interval-clique
rows for flight overlap. This strengthened form is used for all
corrected day-indexed results in \Cref{sec:computational}.

\subsection{Pre-horizon accumulated hours}
\label{subsec:c11_pre_horizon}

The original paper does not account for flying time accumulated before the
planning horizon. In the cumulative-state formulation this history is included
exactly through the initial condition
\begin{equation}\label{eq:c11_pre_horizon}
  h_{j0}=\Phi^A_j,
  \qquad 0\leq\Phi^A_j\leq T^A_{\max},\quad \forall j\in\Pc.
\end{equation}
The condition is required for every aircraft, including $\Phi^A_j=0$; no
separate opening-window big-$M$ inequality is needed.

\subsection{Cumulative-prefix representation of strict interval rows}
\label{subsec:c13_prefix_equivalence}

The strict interval form of C13/C13b also admits a cumulative-prefix
representation. This is an alternative encoding of the interval rows, not the
reset state $h_{jd}$ above. Let $q_{jd}$ be the flight minutes assigned to
aircraft $j$ on day $d$, and define $Q_{j0}=0$ and
\begin{equation}
  Q_{jd}=Q_{j,d-1}+q_{jd},
  \qquad 0\leq Q_{jd}\leq U_j^{\mathrm{fly}},
  \label{eq:c13_prefix_flow}
\end{equation}
where $U_j^{\mathrm{fly}}$ is any valid upper bound on the total assignable
flight minutes for aircraft $j$. For every day pair $d<d'$, telescoping gives
\begin{equation}
  \sum_{i:\,d<\operatorname{depday}(i)\leq d'}
       \tilde t_i x_{ij}=Q_{jd'}-Q_{jd},
  \qquad
  \sum_{i:\,\operatorname{depday}(i)\leq d'}
       \tilde t_i x_{ij}=Q_{jd'}.
  \label{eq:c13_prefix_telescope}
\end{equation}
Consequently, the strict interval rows can be written as
\begin{align}
  Q_{jd'}-Q_{jd}
  &\leq T^A_{\max}+M_{dd'}\sum_{r=d+1}^{d'-1}\mu_{jr},
  \label{eq:c13_prefix_interval}\\
  Q_{jd'}
  &\leq T^A_{\max}-\Phi^A_j
       +M^0\sum_{r=1}^{d'-1}\mu_{jr},
  \label{eq:c13_prefix_opening}
\end{align}
where $\mu_{jr}$ is the A-check reset indicator used by the strict interval
formulation; for A only, $\mu_{jr}=y_{jr}$. The opening row is imposed for
every $d'\in\D$, including $d'=1$, and $M^0$ is a valid bound on the largest
opening-window excess. These rows assume checks after the flights of their
labelled day; they are not sufficient for arbitrary intra-day flexible resets.
Since the prefix flow equations uniquely determine $Q$ from the
assignment variables and the rows above are exactly the telescoped interval
expressions, this substitution preserves the projection onto the common
variables and the bound-preserving LP objective. It does not assert LP
equivalence with the stronger day-indexed reset-state formulation $h_{jd}$.

\subsection{Why the corrected formulations close every gap}
\label{subsec:c13_exactness}

\begin{proposition}[Exactness of the corrected hour rows]
  \label{prop:c13_exactness}
  Let $y$ be binary, $v\ge0$, $0\le\Phi^A_j\le T^A_{\max}$, and let all
  big-$M$ values be valid. Under the day-boundary policy, each of the
  following is satisfied if and only if the physical rule
  $H_{jd}\le T^A_{\max}$ holds for every $j$ and $d$:
  \begin{enumerate}
    \item the state recursion~\eqref{eq:c11_state_lower}--\eqref{eq:c11_state_bound}
      with $h_{j0}=\Phi^A_j$;
    \item the strict interval rows~\eqref{eq:c13_prefix_interval} for all
      $1\le d<d'\le n_d$, together with the opening
      rows~\eqref{eq:c13_prefix_opening} for all $d'\in\D$.
  \end{enumerate}
\end{proposition}

\begin{proof}
  (1) The case analysis following~\eqref{eq:c11_state_bound} shows by
  induction that $h_{jd}$ is forced to equal $H_{jd}$
  in~\eqref{eq:physical_exposure}: without a check on night $d-1$ it adds
  $v_{jd}$, and with one it restarts from $v_{jd}$. Hence the bound
  $h_{jd}\le T^A_{\max}$ is exactly the physical rule.

  (2) Necessity follows from the argument of \Cref{lem:c13_validity}: an
  interval row is active only when $I_{jdd'}=0$, and then
  $L_{jdd'}\le H_{jd'}$. An active opening row for $d'$ has no check before
  $d'$, so its left-hand side plus $\Phi^A_j$ equals $H_{jd'}$. For
  sufficiency, fix $d$. If no check precedes $d$, the opening row for $d$ is
  active and gives $H_{jd}\le T^A_{\max}$. Otherwise let $s$ be the last
  check before $d$. The interval row $(s,d)$ has no interior check, so it is
  active and gives $L_{jsd}=H_{jd}\le T^A_{\max}$.
\end{proof}

The proof also explains why the corrected families are physically more
faithful than C13 and its split. They are anchored to the state that the
limit concerns, the exposure since the last completed check, rather than to
the indicator pattern at the ends of a window. Specifically:
(i) $\Phi^A_j$ enters through $h_{j0}$ or the opening row, closing G1;
(ii) the opening rows relax only on checks strictly before $d'$, closing G2
and the end-day defect of~\eqref{eq:c13b_legacy};
(iii) interval rows exist for every pair $d<d'$, including adjacent days and
gaps of any length, closing G3, G4 and G0;
(iv) a row is active whenever its window contains no interior check,
irrespective of its endpoints, closing G5; and
(v) flights of a check day are counted before that day's reset.
No assumption linking $T^A_{\max}$, $d_{\max}$ and the horizon is needed, and
calendar spacing can be imposed separately without affecting hour coverage.
The state recursion achieves this with $O(n_pn_d)$ continuous variables and
rows, whereas the interval family uses $O(n_pn_d^2)$ rows. Both remain exact
only under the stated day-boundary policy and executable checks; flexible
intra-day resets require the event state of \Cref{sec:event_reformulation}.

\subsection{Integrality constraints}
\begin{align}
  x_{ij}&\in\{0,1\}, && i\in\F,\ j\in\Pc,\\
  z_{ijd}&\in\{0,1\}, && i\in\bigcup_{m\in\MA}F_d^{\downarrow}(m),\ j\in\Pc,\ d\in\D,\\
  y_{jd}&\in\{0,1\}, && j\in\Pc,\ d\in\D.
\end{align}

The variables $h_{jd}$ are continuous and satisfy
$0\leq h_{jd}\leq T^A_{\max}$ as imposed in~\eqref{eq:c11_state_bound}.

There are two independent corrections: C1--C5 supplies a physically executable
assignment, and the cumulative-state recursion enforces the hour threshold
under the stated day-boundary reset policy. The endpoint split is not an
exact substitute. The day-boundary state adds $O(n_p n_d)$ continuous
variables, not $O(n_f n_d)$; a flexible model with flight-event states instead
uses up to $O(n_f n_p)$ hour-state variables. These counts describe different
time resolutions and must not be interchanged.

\subsection{Rigorous feasibility theorem for C1--C13f}
\label{subsec:feasibility_theorem_c1_c13}

We now state the corrected sufficiency result for the hour-only, day-boundary
type-A extension studied here; it is not a claim of equivalence to Khaled's
complete nightly policy.
Consider the day-boundary reference constraint family
\eqref{eq:c0}--\eqref{eq:c4_overlap}, \eqref{eq:c8}--\eqref{eq:c11},
\eqref{eq:c11_state_lower}--\eqref{eq:c11_state_bound}, and the stated
variable domains. Constraint C12 is deliberately omitted in this hour-only
extension: hour spacing is enforced by the cumulative-state rows, but the
original type-A calendar-spacing rule is outside this theorem's assumptions.
The theorem also assumes the overnight executability conditions
in~\eqref{eq:c_night_window} and the accompanying next-departure condition; these
are operational assumptions, not consequences of C8 or C10 alone.

\begin{theorem}[Physical feasibility under day-boundary C1--C13f]
  \label{thm:physical_feasibility_c1_c13}
  Any solution satisfying the above constraints, with $x$, $z$, and $y$
  binary and $h$ continuous, induces for each aircraft $j$ a physically
  executable trajectory over the planning horizon:
  flights assigned to $j$ are temporally non-overlapping, satisfy airport
  continuity and turn time, and admit a day-by-day maintenance schedule
  satisfying capacity and hour-spacing rules.
\end{theorem}

\begin{proof}
  Fix an aircraft $j$ and let
  $S_j=\{i\in\F: x_{ij}=1\}$ ordered by nondecreasing departure time.

  \textbf{(i) Flight assignment and continuity.}
  By Proposition~\ref{prop:corrected_basic_feasibility}, constraints
  C1--C5, including pairwise non-overlap, give each aircraft a continuous,
  turn-time-feasible route. This conclusion does not follow from the balance
  rows~\eqref{eq:c2}--\eqref{eq:c3} alone, as the simultaneous-departure
  counterexample in~\Cref{sec:basic} demonstrates.

  \textbf{(ii) Temporal executability.}
  For any overlapping pair $\{i_1,i_2\}\in\mathcal{O}$,
  \eqref{eq:c4_overlap} enforces $x_{i_1j}+x_{i_2j}\le 1$.
  Hence no two flights assigned to $j$ overlap once turn time is included.
  Combined with (i), this yields a realizable single-aircraft flight timeline.

  \textbf{(iii) Maintenance event consistency.}
  By~\eqref{eq:c9}, a maintenance trigger $z_{ijd}$ can be active only if flight
  $i$ is assigned to $j$. By the aggregation equality~\eqref{eq:c11}, the day
  indicator $y_{jd}$ equals the number of triggers on day $d$; it is binary
  and therefore at most one trigger per aircraft and day is active. The hour
  rows below use $y_{jd}$ only, so every reset is attached to an assigned
  flight and a unique day.

  \textbf{(iv) Maintenance capacity and blocking.}
  By~\eqref{eq:c10}, airport-day capacities are respected. By~\eqref{eq:c8}, if
  maintenance is triggered after flight $i$ on day $d$, then any subsequent
  same-day flight is forbidden for the same aircraft. By
  Assumption~\eqref{eq:c_night_window}, the selected check fits in a staffed
  slot at its arrival airport, and it completes with the required turn time
  before the next assigned departure. Thus the maintenance action is
  operationally executable under that explicit assumption.

  \textbf{(v) Hour/day maintenance admissibility.}
  By~\eqref{eq:c11_state_lower}--\eqref{eq:c11_state_bound}, the cumulative
  flight-time state increases by the flying time assigned each day, resets
  after a maintenance night, and never exceeds $T^A_{\max}$.

  Steps (i)--(v) show that every aircraft has a consistent, non-overlapping
  flight sequence with feasible maintenance placement and resource usage.
  Therefore the global solution is physically feasible.
\end{proof}

%%==========================================================================
\section{Computational Study}
\label{sec:computational}
%%==========================================================================

\subsection{Implementation, instances and protocol}
\label{subsec:setup}

All models are implemented in Python with Pyomo~\citep{pyomo2023}. The
primary exact solves use IBM ILOG CPLEX~12.6~\citep{cplex2023} through Pyomo's
AMPL interface, one thread, presolve enabled and a 180-second wall-clock limit
unless stated otherwise. A complete cross-check with Gurobi~11.0.3
\citep{gurobi2024} (one thread, default presolve, same limit) is reported in
\Cref{subsec:solver_sensitivity}. On Windows, Pyomo's Gurobi shell interface
runs through the Python interpreter bundled with Gurobi, and a model without
a feasible incumbent returns no solution; our drivers preserve that
termination status and do not extract a schedule in that case. The CPLEX
version is old, so absolute timings should not be read as state-of-the-art
solver performance; the Gurobi runs show that objective values do not depend
on this choice.

The instance families are generated deterministically from a seeded
round-trip-rotation generator that certifies each designated rotation with an
independent maintenance counter. This certifies the generated rotations only;
MILP objectives, statuses and schedules come from the solver. Selected
schedules are checked by an independent chronological replay, with coverage
and checked requirements reported for each experiment. In particular, all
loaded uniform-study incumbents are replayed, but the four time-limited
Event-based incumbents in the historical CPLEX study were not replayed.
Where performed, the replay counts a \emph{violation} for every flight after
which the flight-hour counter exceeds the threshold and, when applicable, for
timing, capacity, blocking, coverage or continuity failures. A single overdue
stretch with many flights therefore produces many violations, so counts indicate
the severity of an overdue schedule rather than the number of distinct defects;
the indicator that matters for a given audit is whether its count is zero.
File names, hashes and raw outputs
are listed in \Cref{app:artifacts}. The experiments below are an extended
instance family and not a reproduction of the tables of Khaled et
al.~\citep{khaled2018compact}.

\subsection{Cumulative-flight-time row verification}

The original C13 row, its endpoint-split alternative, and the legacy opening
companion were evaluated by
\texttt{c13\_loophole\_validation.py}, using the hand-built four-day input
\texttt{c13\_loophole\_test} as the base (\Cref{app:artifacts}). The
both-endpoints scenario uses that input unchanged: one aircraft,
$T^A_{\max}=600$ minutes, and 1,200 minutes of flying in $(1,4]$. For the
ending-only opening-gap scenario, the script creates a temporary transformed
copy, changing the day-2 and day-4 flight durations to 590 and 20 minutes,
respectively, for 610 minutes in $(1,4]$. Thus the two rows are not check
patterns on one unchanged duration input. Both runs use $d_{\max}=4$. The
value $M_{14}=1200$ is a valid illustrative bound for these witnesses; the
script evaluates the built model using its configured big-$M$. The table
reports the C13-family result and, for exact-state methods, the corresponding
hour-state result; \Cref{tab:c13_loophole} is not a test of every other model
constraint. With both endpoints checked, C13 is active and the split rows are
relaxed. With only the ending check, C13 is inactive by design; the legacy
opening companion is separately tested for whether it protects pre-check
exposure. The exact state formulations reject both overdue schedules. Thus,
the both-endpoints case is a counterexample to the split, whereas the
one-endpoint case illustrates why a practical model needs a separate
opening-hour condition, not why C13 should be active without two checked
endpoints. The results test
\Cref{subsec:c13_opening_coverage,subsec:c11_corrected}.

\begin{table}[t]
\centering\small
\caption{Row-family numerical verification of conditional C13 and opening-hour
protection. The both-endpoints witness has $T^A_{\max}=600$ and
$\sum_{i\in F_{1,4}}x_{ij}\tilde t_i=1200$; the ending-only opening-gap
scenario uses a temporary transformed input with 610 minutes in the same
window. Both use $d_{\max}=4$. C13 status is row-local; opening status tests
companion C13b and independent replay.}
\label{tab:c13_loophole}
\resizebox{\linewidth}{!}{%
\begin{tabular}{lccccc}
\toprule
Scenario & Khaled C13 & Legacy C13b & Endpoint split & Exact day-state & Exact event-state \\
\midrule
Both endpoints checked & active; rejects & within limit & relaxed; misses & rejects & rejects \\
Ending check only & inactive by design & 0/7 violated; replay fails & active; rejects & rejects & rejects \\
\bottomrule
\end{tabular}}
\end{table}

The historical flexible-start event variant is omitted from this table because it inherits
E9--E13 verbatim from the optimized event model and therefore reproduces its
verdicts on both patterns; the suite run in
\Cref{subsec:event_measurements} confirms this.

\subsection{Practical effect of the non-overlap correction (C5)}
\label{subsec:c5_impact}

The simultaneous-departure gap of \Cref{sec:basic} is a modelling defect, but
whether it changes an optimum depends on the instance. We solved the pure
assignment problem (no maintenance) with Gurobi under the legacy rows C1--C4
and the corrected rows C1--C5 for the six matched instances of
\Cref{subsec:event_measurements} and one smaller instance, and replayed the
C1--C4 incumbent for overlapping flight pairs and continuity breaks. These
timetables contain between 0 and 153 same-airport pairs of flights with
identical departure times, so the corner case is not a rarity. With the
shipped cost matrices, which price each designated rotation aircraft roughly
500 units below every other tail, the C1--C4 optimum is already executable:
objectives coincide in all seven instances and the replay finds no violation.
The defect is thus latent: it matters only if exploiting it is profitable.

To test this, we replaced each cost row by its mean plus 3 per aircraft index,
which leaves only the mild low-index preference of the generator.
\Cref{tab:c5_impact} shows that in four of the seven instances the C1--C4
optimum then places one aircraft on two time-overlapping flights and breaks
route continuity (a ``cloned'' aircraft), and its objective lies strictly
below the true optimum. The understatement is small in relative terms
(0.008--0.035\%), but the plan cannot be flown. In the other three instances
the C1--C4 optimum was already executable. Adding C5 increases the number of
rows by 9--121\% in these instances (clique rows), without a measurable
effect on solution time (all solves took below 3\,s). We conclude that C5
should be included whenever departure times can coincide, because it removes
non-executable schedules. This small-instance result does not establish a
negligible cost at larger scales, nor does it show that C5 changes the optimum
of the shipped instances.

To check that this finding persists at larger scale, we repeated the
experiment on twelve feasible instances with 300--452 flights and 20 or 30
aircraft (three seeds for each of four size cells), again with Gurobi and
both cost variants. All 48 solves were proven optimal. Under the shipped costs,
the C1--C4 and C1--C5 objectives coincide and no violation occurs in any
instance. Under flat costs, every C1--C4 incumbent is non-executable, with
10--105 overlapping pairs and 4--14 continuity breaks, and its objective is
0.04--0.14\% below the true optimum (\Cref{tab:c5_large}); all C1--C5
incumbents are executable. At this scale C5 increases the row count by a factor
of 3--10 and, under flat costs, the solution time by up to a factor of three
(11\,s to 27\,s and 14\,s to 44\,s for $|P|=30$), while under shipped costs the
solution times are unchanged. The cost of C5 is therefore model size, and at
this scale also some solution time when the cost structure leaves room for
exploiting the defect, not a loss of tractability.

We extended the check to twelve further instances with 600--752 flights and
40 or 50 aircraft (three seeds per cell, 200\,s limit per solve). Under the
shipped costs all 24 solves were optimal, the C1--C4 and C1--C5 objectives
coincide, and the replay finds no violation; the C1--C5 model has 7--19 times
as many rows (for example 32\,600 versus 485\,240 for $|P|=40$, $H=16$) but
the same solution time (24--65\,s). Under flat costs the flat-cost problem
becomes much harder, and only 7 of the 24 solves finished within the limit
(\Cref{tab:c5_p4050}). Every C1--C4 incumbent returned, optimal or not,
contains overlapping flights and continuity breaks (9--24 breaks per
instance on average), and no C1--C5 incumbent does. Because most flat-cost
solves stopped at the limit with large gaps, no cell allows a paired
comparison of proven optima, and we make no claim about objective differences
at this scale.

\begin{table}[t]
\centering\small
\caption{Larger instances under flat costs (Gurobi, 200\,s limit, means over
three seeds). ``Opt'' is the number of seeds solved to optimality. Overlaps and
breaks are counted on the returned incumbent, optimal or not.}
\label{tab:c5_p4050}
\begin{tabular}{rrr rrrr rrrr}
\toprule
 & & & \multicolumn{4}{c}{C1--C4} & \multicolumn{4}{c}{C1--C5} \\
\cmidrule(lr){4-7}\cmidrule(lr){8-11}
$|P|$ & $H$ & $F$ & Opt & Time (s) & Overl. & Breaks & Opt & Time (s) & Overl. & Breaks \\
\midrule
40 & 16 & 600 & 0/3 & 203 & 190 & 19 & 3/3 & 131 & 0 & 0 \\
40 & 31 & 600 & 3/3 & 62  & 45  & 9  & 0/3 & 203 & 0 & 0 \\
50 & 16 & 750 & 1/3 & 197 & 300 & 24 & 0/3 & 207 & 0 & 0 \\
50 & 31 & 752 & 0/3 & 394 & 74  & 12 & 0/3 & 208 & 0 & 0 \\
\bottomrule
\end{tabular}
\end{table}

\begin{table}[t]
\centering\small
\caption{Effect of C5 on larger instances under flat costs (Gurobi, optimal in
all runs; means over three seeds). Overlaps and breaks are counted on the
replayed C1--C4 incumbent; C1--C5 incumbents have none. Under shipped costs the
objectives coincide and no violation occurs.}
\label{tab:c5_large}
\begin{tabular}{rrr rr r rr rr}
\toprule
$|P|$ & $H$ & $F$ & Obj.\ C1--C4 & Obj.\ C1--C5 & Gap (\%) & Overlaps & Breaks & Rows C1--C4 & Rows C1--C5 \\
\midrule
20 & 16 & 300 & 290252 & 290522 & 0.093 & 45  & 9  & 10300 & 65620 \\
20 & 31 & 300 & 290365 & 290485 & 0.041 & 10  & 4  & 13900 & 41020 \\
30 & 16 & 450 & 445636 & 446266 & 0.141 & 105 & 14 & 19950 & 209430 \\
30 & 31 & 452 & 447978 & 448314 & 0.075 & 28  & 7  & 25412 & 120032 \\
\bottomrule
\end{tabular}
\end{table}

\begin{table}[t]
\centering\small
\caption{Effect of C5 under flat costs (Gurobi, optimal in all runs). ``Equal''
is the number of same-airport flight pairs with identical departure time.
Overlaps and breaks are counted on the replayed C1--C4 incumbent; C1--C5
incumbents have none. Shipped costs give identical objectives and no
violations on all seven instances.}
\label{tab:c5_impact}
\begin{tabular}{rrrr rrr rr}
\toprule
$|P|$ & $H$ & $F$ & Equal & Obj.\ C1--C4 & Obj.\ C1--C5 & Gap (\%) & Overlaps & Breaks \\
\midrule
10 & 7  & 50  & 5   & 46276  & 46276  & 0     & 0 & 0 \\
5  & 30 & 74  & 8   & 64445  & 64457  & 0.019 & 1 & 1 \\
5  & 48 & 80  & 0   & 69699  & 69699  & 0     & 0 & 0 \\
10 & 30 & 146 & 59  & 135511 & 135547 & 0.027 & 3 & 2 \\
10 & 49 & 160 & 32  & 148578 & 148578 & 0     & 0 & 0 \\
15 & 30 & 218 & 153 & 207587 & 207659 & 0.035 & 6 & 3 \\
15 & 49 & 240 & 96  & 228644 & 228662 & 0.008 & 1 & 1 \\
\bottomrule
\end{tabular}
\end{table}

\subsection{Uniform flexible-start study}
\label{subsec:uniform_results}
The main comparison uses the common policy of \Cref{subsec:uniform_flexible}.
Both encodings use CPLEX 12.6.0.0, one thread, presolve on, and a 180-second limit.
They share the hour engine, so the comparison measures the maintenance encoding and not independent physical models.
All loaded incumbents, including time-limited ones, are checked by chronological replay.
An additional audit reconstructs cost, hours, timing, routes, actual start-day counts and bay occupancy
directly from the saved decisions and the input data, without the model's state values.
\Cref{tab:uniform_flexible} reports the six matched instances.
\begin{table}[t]\centering\scriptsize
\caption{Uniform flexible-start encodings. Wall time is measured around the solver call.
``Limit'' means the 180-second limit was reached with an incumbent; such an incumbent is not a proven optimum.}
\label{tab:uniform_flexible}\resizebox{\linewidth}{!}{%
\begin{tabular}{rrl lrrrrl}\toprule
$|P|$ & $H$ & Encoding & Status & Objective & Wall (s) & Vars & Rows & Replay \\\midrule
5 & 30 & Day & optimal & 37583 & 0.84 & 1548 & 5283 & PASS \\
5 & 30 & Event & optimal & 37583 & 0.96 & 1733 & 5468 & PASS \\
5 & 48 & Day & optimal & 40707 & 0.73 & 1560 & 5360 & PASS \\
5 & 48 & Event & optimal & 40707 & 0.70 & 1760 & 5560 & PASS \\
10 & 30 & Day & optimal & 74150 & 51.38 & 6188 & 21220 & PASS \\
10 & 30 & Event & optimal & 74150 & 27.17 & 6918 & 21950 & PASS \\
10 & 49 & Day & optimal & 81391 & 3.15 & 6463 & 21959 & PASS \\
10 & 49 & Event & optimal & 81391 & 3.00 & 7263 & 22759 & PASS \\
15 & 30 & Day & optimal & 110749 & 157.11 & 13900 & 49043 & PASS \\
15 & 30 & Event & limit & 110749 & 182.39 & 15535 & 50678 & PASS \\
15 & 49 & Day & limit & 122169 & 182.57 & 14607 & 50211 & PASS \\
15 & 49 & Event & limit & 122169 & 182.41 & 16407 & 52011 & PASS \\
\bottomrule\end{tabular}}\end{table}
Of 12 runs, 9 terminate optimally and 12 return loaded incumbents; all 12 pass both replay checks.
Solver bounds and node counts that the solver interface did not return are left blank in the raw data, not inferred from incumbent values.
A Gurobi rerun (\Cref{subsec:solver_sensitivity}) returns the same objective in all twelve runs and stops at its limit on the same three, with residual gaps of 0.7--1.3\%. For $P=15$, $H=30$, the time-limited Event incumbent matches the Day encoding's solver-certified optimal objective; by the equivalence in \Cref{prop:uniform_equivalence}, this establishes the same optimum for the Event encoding, subject to the solver's optimality tolerance, although its own run did not certify it. The two time-limited incumbents at $P=15$, $H=49$ remain best-known values without a proven optimum.
The timings are descriptive workstation measurements, not isolated repeated-runtime trials or evidence of general speedup.
\begin{table}[t]\centering\scriptsize
\caption{Repaired day-boundary ablation, including the first-day opening row. This is not the arbitrary flexible-start model.}
\label{tab:repaired_prefix}\resizebox{\linewidth}{!}{%
\begin{tabular}{rr rrr r ll}\toprule
$|P|$ & $H$ & State LP & Interval LP & Prefix LP & Common MIP & Interval vars/rows & Prefix vars/rows \\\midrule
10 & 30 & 72050.00 & 67850.11 & 67850.11 & 74150 & 3120/12395 & 3440/12705 \\
10 & 49 & 79057.67 & 74391.13 & 74391.13 & 81391 & 3900/20798 & 4410/21298 \\
5 & 30 & 36516.33 & 34383.06 & 34383.06 & 37583 & 1020/4498 & 1180/4653 \\
5 & 48 & 39557.00 & 37207.06 & 37207.06 & 40707 & 1335/8625 & 1585/8870 \\
\bottomrule\end{tabular}}\end{table}
All 12 repaired LP and all 12 integer runs terminate optimally. Saved integer schedules pass the day-boundary hour replay.
The interval/prefix LP equality follows from telescoping with identical bounds and opening domains; no general LP-strength ordering with the day-state formulation is claimed.
\Cref{tab:repaired_prefix} is a day-boundary ablation and not a result for the flexible-start encodings.
\Cref{fig:uniform_variables,fig:uniform_wall} show model size and solve time at $H=30$.
\begin{figure}[t]\centering
\includegraphics[width=0.86\linewidth]{figures/uniform_variables_h30_20261021.png}
\caption{Uniform flexible model size at observed $H=30$, including the P=15 models.}
\label{fig:uniform_variables}
\end{figure}
\begin{figure}[t]\centering
\includegraphics[width=0.86\linewidth]{figures/uniform_wall_h30_20261021.png}
\caption{Uniform flexible solve time at $H=30$ on common optimal cells only (P=5 and P=10).}
\label{fig:uniform_wall}
\end{figure}
\Cref{tab:uniform_night} reports the fixed 22:00 restriction.
\begin{table}[t]\centering\scriptsize
\caption{Fixed 22:00 restriction on the four P=5/10 cases, with the same 180-second limit.}
\label{tab:uniform_night}\begin{tabular}{rrllrrl}\toprule
$|P|$ & $H$ & Encoding & Status & Objective & Wall (s) & Replay \\\midrule
5 & 30 & Day & optimal & 37583 & 0.78 & PASS \\
5 & 30 & Event & optimal & 37583 & 0.78 & PASS \\
5 & 48 & Day & optimal & 40707 & 0.69 & PASS \\
5 & 48 & Event & optimal & 40707 & 0.69 & PASS \\
10 & 30 & Day & optimal & 74150 & 12.76 & PASS \\
10 & 30 & Event & optimal & 74150 & 12.75 & PASS \\
10 & 49 & Day & optimal & 81391 & 2.15 & PASS \\
10 & 49 & Event & optimal & 81391 & 2.18 & PASS \\
\bottomrule\end{tabular}\end{table}
The nightly model is a restricted flexible model, not Khaled's daily-quota/calendar-spacing model.
\Cref{tab:uniform_seeds} gives a three-seed CPLEX replication; the ten-seed Gurobi replication of \Cref{subsec:ten_seed} supersedes it for timing comparisons.
\begin{table}[t]\centering\scriptsize
\caption{Uniform flexible three-seed subset: wall time as mean $\pm$ sample standard deviation, optimal runs only.}
\label{tab:uniform_seeds}\begin{tabular}{rrllr}\toprule
$|P|$ & $H$ & Encoding & Optimal / runs & Wall (s) \\\midrule
5 & 30 & Day & 3/3 & $0.874\pm0.027$ \\
5 & 30 & Event & 3/3 & $0.892\pm0.084$ \\
5 & 48 & Day & 3/3 & $0.714\pm0.005$ \\
5 & 48 & Event & 3/3 & $0.706\pm0.009$ \\
10 & 30 & Day & 3/3 & $72.943\pm52.382$ \\
10 & 30 & Event & 3/3 & $22.516\pm5.142$ \\
10 & 49 & Day & 3/3 & $2.705\pm0.072$ \\
10 & 49 & Event & 3/3 & $2.893\pm0.076$ \\
\bottomrule\end{tabular}\end{table}
Three seeds test limited instance variation; seed zero overlaps the six-case study and is not an additional independent sample.
Raw results, saved schedules, source archives and hashed metadata are listed in \Cref{app:artifacts}.


\subsection{Historical five-method type-A comparison}
\label{subsec:event_measurements}

The ordered implementation removes explicit successor variables and propagates
an aircraft-local prefix state through flights sorted by departure time. A
maintenance event is attached to its triggering flight, and the pre-check row
prevents an event from repairing an already overdue arrival. The five methods
evaluated in the matched A-only experiment are the original Khaled row,
the endpoint split, the strengthened corrected legacy rows, the optimized
ordered event-based formulation, and its flexible-start variant. A result is
described as objective-matched only when both compared methods are optimal.
The legacy methods use a zero-day trigger-to-check deferral cap in the matched
rerun; the optimized event model starts maintenance after the triggering
flight, while the flexible event model permits a one-day (1440-minute) offset.
The day-indexed and event-based timing rules remain different, so this setting
improves comparability but does not make the formulations identical. This
subsection and the following ones therefore use a different maintenance
policy from the uniform study above; their values are not results for the
uniform encodings. In particular, the day-boundary interval/prefix ablation
below is not evidence for the repaired formulation of
\Cref{tab:repaired_prefix}.

\subsubsection{Longer horizons and fleet sizes}

The five methods were run on six matched A-only instances generated instances (\Cref{app:artifacts}). Fleet sizes are
$|P|\in\{5,10,15\}$, observed horizons are $H\in\{30,48,49\}$, and the
instances contain 74--240 flights with up to eight seeded round-trip
rotations per aircraft. Each rotation is checked by the generator's
independent maintenance replay; this validates the designated generated
rotations, not every schedule returned by a MILP. All five formulations use
the same JSON instance within each cell. CPLEX~12.6 was run with one thread,
presolve enabled, and a 180-second limit. Legacy methods were run with
\texttt{max\_hour\_check\_deferral\_days=0}; event-flex retains its 1440-minute
deferral window. Results, input hashes and solver configuration are archived
with the paper (\Cref{app:artifacts}). This is an extended-family experiment, not a reproduction of
Khaled et al.'s Table~6.

Table~\ref{tab:longflight_results} reports solver status, incumbent objective,
timing, model size, and independent replay status for all 30 performance runs.
Both Event-based methods reached the time limit at P=15 for H=30 and H=49;
their incumbents are not proven optima and were not independently replayed.

\begin{table}[p]
\centering\scriptsize
\caption{Five-method A-only results on the longer-horizon, higher-flight-count
family with zero-day legacy check deferral. ``Replay'' is the independent
flight-hour audit (PASS or violation count); time-limited rows are not
optimality certificates. CPU and wall times are in seconds.}
\label{tab:longflight_results}
\resizebox{\linewidth}{!}{%
\begin{tabular}{rrrlllr rrrr}
\toprule
$|P|$ & $H$ & $F$ & Method & Status & Replay & Objective & CPU & Wall & Vars & Cons. \\
\midrule
5 & 30 & 74 & Legacy paper & optimal & 51 & 34883 & 0.52 & 1.02 & 1020 & 4348 \\
5 & 30 & 74 & Legacy split & optimal & 17 & 36083 & 1.01 & 1.71 & 1020 & 6523 \\
5 & 30 & 74 & Legacy corrected & optimal & PASS & 37583 & 0.32 & 0.67 & 1175 & 2338 \\
5 & 30 & 74 & Event optimized & optimal & PASS & 37583 & 0.33 & 0.67 & 925 & 3007 \\
5 & 30 & 74 & Event flex & optimal & PASS & 37583 & 0.41 & 0.88 & 1363 & 4393 \\
5 & 48 & 80 & Legacy paper & optimal & 60 & 37707 & 0.49 & 1.36 & 1335 & 8385 \\
5 & 48 & 80 & Legacy split & optimal & 15 & 39207 & 0.92 & 2.24 & 1335 & 14025 \\
5 & 48 & 80 & Legacy corrected & optimal & PASS & 40707 & 0.28 & 0.67 & 1580 & 3310 \\
5 & 48 & 80 & Event optimized & optimal & PASS & 40707 & 0.30 & 0.67 & 1000 & 3240 \\
5 & 48 & 80 & Event flex & optimal & PASS & 40707 & 0.29 & 0.68 & 1360 & 4320 \\
10 & 30 & 146 & Legacy paper & optimal & 106 & 68850 & 0.94 & 1.96 & 3120 & 12095 \\
10 & 30 & 146 & Legacy split & optimal & 30 & 71150 & 2.56 & 4.05 & 3120 & 16445 \\
10 & 30 & 146 & Legacy corrected & optimal & PASS & 74150 & 1.55 & 2.13 & 3430 & 5725 \\
10 & 30 & 146 & Event optimized & optimal & PASS & 74150 & 12.32 & 12.93 & 3650 & 11425 \\
10 & 30 & 146 & Event flex & optimal & PASS & 74150 & 10.38 & 11.15 & 5458 & 18000 \\
10 & 49 & 160 & Legacy paper & optimal & 120 & 75391 & 3.29 & 5.36 & 3900 & 20308 \\
10 & 49 & 160 & Legacy split & optimal & 30 & 78391 & 10.12 & 13.62 & 3900 & 32068 \\
10 & 49 & 160 & Legacy corrected & optimal & PASS & 81391 & 0.77 & 1.53 & 4400 & 8048 \\
10 & 49 & 160 & Event optimized & optimal & PASS & 81391 & 0.98 & 2.04 & 4000 & 12672 \\
10 & 49 & 160 & Event flex & optimal & PASS & 81391 & 1.57 & 2.37 & 5663 & 18269 \\
15 & 30 & 218 & Legacy paper & optimal & 158 & 102849 & 1.72 & 3.73 & 6300 & 25822 \\
15 & 30 & 218 & Legacy split & optimal & 46 & 106249 & 5.68 & 8.62 & 6300 & 32347 \\
15 & 30 & 218 & Legacy corrected & optimal & PASS & 110749 & 12.49 & 13.48 & 6765 & 10192 \\
15 & 30 & 218 & Event optimized & maxIterations & -- & 110749 & 180.41 & 181.50 & 8175 & 25222 \\
15 & 30 & 218 & Event flex & maxIterations & -- & 110749 & 180.50 & 182.08 & 12265 & 42053 \\
15 & 49 & 240 & Legacy paper & optimal & 168 & 113169 & 4.25 & 8.57 & 7650 & 36633 \\
15 & 49 & 240 & Legacy split & optimal & 45 & 117669 & 15.42 & 22.45 & 7650 & 54273 \\
15 & 49 & 240 & Legacy corrected & optimal & PASS & 122169 & 24.10 & 25.59 & 8400 & 13863 \\
15 & 49 & 240 & Event optimized & maxIterations & -- & 122169 & 180.45 & 181.75 & 9000 & 28039 \\
15 & 49 & 240 & Event flex & maxIterations & -- & 122169 & 180.52 & 182.35 & 12807 & 42276 \\
\bottomrule
\end{tabular}
}%
\end{table}

With zero-day legacy deferral, the strengthened legacy and both Event-based
formulations have identical optimal objectives on the four cells where all
three methods prove optimality. On the two P=15 cells, the Event-based methods
return the same incumbent as strengthened legacy but hit the time limit, so
that agreement is not an optimality certificate. The original-paper and
endpoint-split models return lower objectives, but their CPLEX-selected
schedules exceed the A-hour threshold (663 and 183 replay violations across
the six cases, respectively); those objectives are not feasible-quality
comparisons. The solver sensitivity of these counts is examined below.
The optimized Event-based model has fewer variables than strengthened legacy
on three of six cells; averaged over all six, its variable count is slightly
higher. Neither Event-based formulation consistently reduces constraints or
timing. These results are limited to this generated family and solver setup.

\subsubsection{Cross-solver check with Gurobi}
\label{subsec:solver_sensitivity}
To assess whether the replay counts depend on the optimizer-selected
schedule, we reran the same six hashed instances with Gurobi~11.0.3, one
thread, presolve enabled, and the same 180-second limit. All 30 runs returned
incumbents and terminated with an optimal status under the solver's default
MIP-gap tolerance; the event-flex run at $P=15$, $H=49$ had a residual gap of
0.0041\%. A separate zero-gap run closed this gap at the same objective.
Objectives for each formulation-instance pair matched the CPLEX results.
Table~\ref{tab:replay_solver_sensitivity} compares the independent
flight-hour replay counts for the two formulations whose returned schedules
failed that audit.

\begin{table}[t]
\centering\small
\caption{Flight-hour replay violations in solver-selected optimal schedules
on the six matched A-only instances. Counts characterize the returned
schedules, not the formulations alone.}
\label{tab:replay_solver_sensitivity}
\begin{tabular}{rrrrrr}
\toprule
 & & \multicolumn{2}{c}{Original paper C13} &
 \multicolumn{2}{c}{Endpoint split} \\
\cmidrule(lr){3-4}\cmidrule(lr){5-6}
$|P|$ & $H$ & CPLEX & Gurobi & CPLEX & Gurobi \\
\midrule
5  & 30 & 51  & 51  & 17 & 15 \\
5  & 48 & 60  & 60  & 15 & 15 \\
10 & 30 & 106 & 106 & 30 & 30 \\
10 & 49 & 120 & 117 & 30 & 37 \\
15 & 30 & 158 & 158 & 46 & 47 \\
15 & 49 & 168 & 180 & 45 & 56 \\
\midrule
Total & & 663 & 672 & 183 & 200 \\
\bottomrule
\end{tabular}
\end{table}

The totals differ between solvers despite matching objective values,
demonstrating that multiple optimal schedules can yield different replay
counts. In both runs, however, every original-C13 and endpoint-split
incumbent failed the independent hour audit, while all strengthened-legacy
and Event-based incumbents passed. The Gurobi run also proved optimality for
the four P=15 Event-based cases that reached the CPLEX time limit; their
schedules passed replay. Raw results, schedules and provenance for these runs
are listed in \Cref{app:artifacts}.

\subsubsection{Ten-seed Gurobi replication of the uniform encodings}
\label{subsec:ten_seed}
To obtain a less anecdotal timing comparison of the two uniform encodings, we
generated ten deterministic replicates (indices 0--9; indices 0--2 coincide
with the earlier three-seed set) for each of the four P=5/10 cells and solved
both encodings with Gurobi (one thread, 180\,s limit). All 80 runs terminated
optimally and passed both replay audits; the two encodings returned the same
optimal objective on every instance. \Cref{tab:ten_seed} summarizes the wall
times. The event-indexed encoding was faster in 9, 8, 7 and 6 of the ten
replicates in the four cells; a two-sided Wilcoxon signed-rank test gives
$p<0.05$ in two cells ($p=0.010$ and $0.027$, not corrected for multiple
comparisons), and the geometric-mean time ratios (event over day) are 0.87,
0.96, 0.60 and 0.96. The effect sizes are modest, the
absolute times are small, and the earlier CPLEX three-seed subset shows a
different sign at $|P|=10$, $H=49$. We therefore regard the encodings as
computationally comparable on this family, with a possible advantage of the
event-indexed one at the hardest cell, and make no solver-independent claim.

\begin{table}[t]
\centering\small
\caption{Ten-seed Gurobi replication of the uniform flexible encodings (all
runs optimal, all replay checks passed). Wall time in seconds as mean
$\pm$ sample standard deviation over ten instances; ``Faster'' is the number
of instances on which the event-indexed encoding was faster.}
\label{tab:ten_seed}
\begin{tabular}{rr r rr r}
\toprule
$|P|$ & $H$ & Mean objective & Day & Event & Faster \\
\midrule
5  & 30 & 37\,619 & $0.61\pm0.08$ & $0.53\pm0.04$ & 9/10 \\
5  & 48 & 40\,705 & $0.49\pm0.04$ & $0.47\pm0.06$ & 8/10 \\
10 & 30 & 74\,190 & $15.06\pm6.33$ & $8.58\pm1.64$ & 7/10 \\
10 & 49 & 81\,404 & $3.07\pm0.62$ & $2.95\pm0.55$ & 6/10 \\
\bottomrule
\end{tabular}
\end{table}

The solve times of the P=10 cells over the observed horizon are shown in
\Cref{fig:ordered_a_wall_horizon,fig:ordered_a_cpu_horizon}.

\begin{figure}[t]
\centering
\includegraphics[width=0.86\linewidth]{figures/ordered_a_wall_vs_horizon_extended_p10.png}
\caption{A-only solve wall time versus observed horizon for $|P|=10$. Only
paired cells where all five methods are optimal are plotted.}
\label{fig:ordered_a_wall_horizon}
\end{figure}

\begin{figure}[t]
\centering
\includegraphics[width=0.86\linewidth]{figures/ordered_a_cpu_vs_horizon_extended_p10.png}
\caption{A-only solve CPU time versus observed horizon for $|P|=10$. Only
paired cells where all five methods are optimal are plotted.}
\label{fig:ordered_a_cpu_horizon}
\end{figure}

\subsubsection{Fleet-size scaling at \texorpdfstring{$H=30$}{H=30}}

The fleet-size figures (\Cref{fig:a_variables_fleet,fig:a_constraints_fleet,fig:a_cpu_fleet,fig:a_wall_fleet}) use the complete matched set at observed horizon
$H=30$ for $|P|\in\{5,10,15\}$. Model-size figures include all successfully
built models; timing figures include only cells where all five methods are
optimal (P=5 and P=10).

\begin{figure}[t]
\centering
\includegraphics[width=0.86\linewidth]{figures/a_variables_vs_fleet_extended_h30.png}
\caption{Active Pyomo variables versus fleet size at observed horizon $H=30$.}
\label{fig:a_variables_fleet}
\end{figure}

\begin{figure}[t]
\centering
\includegraphics[width=0.86\linewidth]{figures/a_constraints_vs_fleet_extended_h30.png}
\caption{Active Pyomo constraints versus fleet size at observed horizon $H=30$.}
\label{fig:a_constraints_fleet}
\end{figure}

\begin{figure}[t]
\centering
\includegraphics[width=0.86\linewidth]{figures/a_cpu_vs_fleet_extended_h30.png}
\caption{A-only solve CPU time versus fleet size at observed horizon $H=30$.
Only paired fleet sizes where all five methods are optimal are plotted.}
\label{fig:a_cpu_fleet}
\end{figure}

\begin{figure}[t]
\centering
\includegraphics[width=0.86\linewidth]{figures/a_wall_vs_fleet_extended_h30.png}
\caption{A-only solve wall time versus fleet size at observed horizon $H=30$.
Only paired fleet sizes where all five methods are optimal are plotted.}
\label{fig:a_wall_fleet}
\end{figure}

\subsubsection{Summary of the five-method results}

All 30 solver runs returned incumbents and 26 terminated optimally. The four
time-limited Event-based runs are the P=15 cases at H=30 and H=49; their
incumbents are not optimality certificates and are excluded from paired timing
averages. Independent replay reports zero violations for all six strengthened
legacy schedules and the four optimal schedules from each Event-based method.
The original-paper and endpoint-split schedules fail replay in all six cells,
consistent with the C13 counterexamples above. Time-limited Event-based
schedules were not replayed.

\subsubsection{Aggregate model size and optimal-run timing}

Table~\ref{tab:longflight_aggregates} averages model counts over all six
performance cases. Timing averages use the same four P=5/10 cells, where all
five methods terminated optimally. This paired subset avoids comparing means
computed from different sets of solved instances.

\begin{table}[t]
\centering\scriptsize
\caption{Mean active scalar model size over six generated cases; mean build,
CPU, and wall times over the four common optimal cells. Times are in seconds.}
\label{tab:longflight_aggregates}
\resizebox{\linewidth}{!}{%
\begin{tabular}{lrrrrrr}
\toprule
Method & Optimal/6 & Mean vars/6 & Mean cons./6 & Mean build/4 & Mean CPU/4 & Mean wall/4 \\
\midrule
Legacy paper & 6/6 & 3887.5 & 17931.8 & 1.327 & 1.310 & 2.426 \\
Legacy split & 6/6 & 3887.5 & 25946.8 & 1.412 & 3.654 & 5.404 \\
Legacy corrected & 6/6 & 4291.7 & 7246.0 & 0.181 & 0.730 & 1.248 \\
Event optimized & 4/6 & 4458.3 & 13934.2 & 0.325 & 3.486 & 4.080 \\
Event flex & 4/6 & 6486.0 & 21551.8 & 0.334 & 3.161 & 3.766 \\
\bottomrule
\end{tabular}}
\end{table}

On the four paired optimal cells, the corrected legacy method has the lowest
mean CPU and wall time; the ordered formulations do not show a timing advantage
on this subset. Average model sizes are descriptive and depend on the chosen
domain and maintenance-timing semantics. The original and split rows have
lower objectives but violate the independent hour replay, so they are not
feasible performance baselines.

\subsubsection{Cumulative-prefix LP and MILP ablation}
\label{subsec:hour_prefix_ablation}

\Cref{tab:hour_prefix_ablation} compares the strict interval formulation with an opt-in cumulative
daily-prefix representation on the four P=5/10 instances. Both use the same
strict C13/C13b conditions, zero-day legacy deferral, and bound-preserving LP
relaxation. All interval and prefix LPs and MILPs terminated optimally on these
four cases. The prefix projection reproduces the interval LP objective to
numerical precision, as expected from the telescoping reformulation in
\Cref{subsec:c13_prefix_equivalence}, and the two formulations return identical integer
objectives. In this four-instance ablation, the exact day-state model has a
stronger LP bound on each instance; this empirical comparison does not imply
a general LP-strength ordering. Prefix-versus-interval equivalence should not
be misread as equivalence to the day-state relaxation. The P=15 LP ablation is excluded:
several solver runs reached their iteration limit and returned no reliable LP
objective.

\begin{table}[t]
\centering\scriptsize
\caption{A-only cumulative-prefix ablation on four tractable instances. LP
columns report optimal lower bounds; the MIP objective is shared by exact
day-state, strict interval, and prefix models. Prefix model size and MIP wall
time are shown for comparison.}
\label{tab:hour_prefix_ablation}
\resizebox{\linewidth}{!}{%
\begin{tabular}{rrr rrr r r r r}
\toprule
$|P|$ & $H$ & $F$ & Exact-state LP & Interval LP & Prefix LP & MIP objective & Interval vars/rows & Prefix vars/rows & Prefix wall (s) \\
\midrule
5 & 30 & 74 & 36516.33 & 34383.06 & 34383.06 & 37583 & 1020/4493 & 1180/4648 & 0.91 \\
5 & 48 & 80 & 39557.00 & 37207.06 & 37207.06 & 40707 & 1335/8620 & 1585/8865 & 1.05 \\
10 & 30 & 146 & 72050.00 & 67850.11 & 67850.11 & 74150 & 3120/12385 & 3440/12695 & 8.24 \\
10 & 49 & 160 & 79057.67 & 74391.13 & 74391.13 & 81391 & 3900/20788 & 4410/21288 & 4.61 \\
\bottomrule
\end{tabular}
}%
\end{table}

The prefix representation adds continuous states and flow rows, and it does
not improve solve time uniformly: it is faster than interval on two cells and
slower on two. This is a formulation-equivalence and size observation for the
tested A-only subset, not evidence of a general performance advantage. Raw
LP/MILP rows and run metadata are listed in \Cref{app:artifacts}.

\subsubsection{Three-seed A-only replication subset}

To measure instance-to-instance variation, we generated three deterministic
replicates for each of the four P=5/10, H=30/48/49 cells (\Cref{app:artifacts}). All 60 method-instance runs
terminated optimally. The strengthened legacy and both Event-based objectives
matched on every replicate. Table~\ref{tab:a_only_multiseed} reports the mean
common objective and wall time as mean $\pm$ sample standard deviation. These
three replicates are a limited robustness check, not a basis for significance
tests or a general runtime ranking.

\begin{table}[t]
\centering\scriptsize
\caption{Three-seed A-only subset. Each cell has three generated instances;
all 60 runs are optimal. Objective means agree across the three corrected
methods. Wall time is mean $\pm$ sample standard deviation in seconds.}
\label{tab:a_only_multiseed}
\begin{tabular}{rrr r rrr}
\toprule
$|P|$ & $H$ & $F$ & Mean objective & Strengthened legacy & Event optimized & Event flex \\
\midrule
5 & 30 & 74 & 37603.7 & $0.689\pm0.008$ & $0.700\pm0.011$ & $0.812\pm0.040$ \\
5 & 48 & 80 & 40706.3 & $0.695\pm0.004$ & $0.642\pm0.007$ & $0.679\pm0.008$ \\
10 & 30 & 146 & 74219.3 & $2.064\pm0.131$ & $13.099\pm7.063$ & $12.975\pm2.532$ \\
10 & 49 & 160 & 81405.3 & $1.527\pm0.055$ & $1.843\pm0.265$ & $2.395\pm0.050$ \\
\bottomrule
\end{tabular}
\end{table}

The P=10, H=30 Event-based wall-time standard deviations are large relative
to their means, illustrating why the single-instance scale-up timings should
not be interpreted as stable method rankings. All five methods pass the
objective/status checks in this subset, but the independent replay still
rejects the paper-C13 and endpoint-split schedules; their lower objectives
must not be treated as feasible comparisons. Raw rows and provenance are listed in
\Cref{app:artifacts}.

\subsubsection{Conditional removal of unnecessary A checks}

We also tested the sufficient pruning condition that removes optional A-check
triggers when a tail's initial exposure plus all assignable flight hours is
below the A threshold. On a modified six-flight A-only control, each tail had
at most 20 assignable flight hours and zero initial exposure, versus a
100{,}000-minute threshold (1{,}666.7 hours). The unpruned and pruned models
both solved optimally at objective 600 and passed independent replay. Pruning
removed all 12 trigger variables and reduced the model from 88 to 76 variables
and from 187 to 172 constraints. This is a construction/model-size check on a
deliberately qualifying instance; it does not establish a runtime improvement
on operationally binding cases. Table~\ref{tab:a_check_pruning} summarizes the
run; hashes of the run and the transformed input are archived
(\Cref{app:artifacts}).

\begin{table}[t]
\centering\small
\caption{Conditional A-check pruning on a constructed qualifying instance.}
\label{tab:a_check_pruning}
\begin{tabular}{lrrrrrrl}
\toprule
Variant & Status & Objective & $z$ vars & Vars & Cons. & Wall (s) & Replay \\
\midrule
Baseline & optimal & 600.0 & 12 & 88 & 187 & 0.544 & PASS \\
Pruned & optimal & 600.0 & 0 & 76 & 172 & 0.513 & PASS \\
\bottomrule
\end{tabular}
\end{table}

\subsubsection{Independent schedule replay and scope}

The generator replay certifies the designated rotations; the separate
post-solve replay audits optimizer-selected schedules. For legacy models it
resets the hour counter after the selected maintenance day; for ordered event
models it resets after the selected trigger flight and independently checks
event timing, flight blocking, and station capacity. This protocol detects
violations in the paper and split C13 variants and returns zero violations for
the strengthened exact-state and optimal Event-based schedules in this suite.
It does not certify the four Event-based incumbents stopped by the time limit.

\subsubsection{Event-model comparison}

Optimized and flexible Event-based objectives match the strengthened legacy
optimum on the four jointly optimal P=5/10 cells. At P=15, both event methods
return the same incumbent as strengthened legacy, but neither proves
optimality at either horizon. No objective-equivalence claim is made beyond
these matched cases.

\subsubsection{Flexible-start Event-based variant}

The flexible-start implementation uses the integer-grid E14--E21 formulation
with a one-day deferral window, as described in
\Cref{subsec:flex_complexity}. Its model has more variables and constraints
than the immediate-start event model in each generated performance cell. The
objective comparison above is limited to solver results and does not replace
independent schedule validation.





\subsection{Scale check with 40 and 50 aircraft}
\label{subsec:p4050}

To probe Khaled-size fleets, we solved twelve type-A instances with 40 or 50
aircraft and 600--752 flights (horizons 15 and 30 days, three seeds per cell)
with Gurobi, one thread and a 200\,s limit per solve. Replay rows count
violations in the returned incumbent. The published row alone
(\emph{legacy paper}) was solved to optimality in all 12 cases but violates
the flight-hour limit in every one (mean 480 violations); the endpoint split
was optimal in 6 of 12 cases and its incumbents violate the limit in all 12;
the exact-state row with a pre-check guard was optimal in 6 of 12 and
otherwise stopped with a mean gap of 0.86\%, and its incumbents were
violation-free throughout; the event-indexed model returned a violation-free
incumbent in all 12 cases but proved none optimal within the limit (mean gap
1.5\%). Solution times of the exact formulations therefore grow faster with
fleet size than those of the unsafe rows, which solve quickly because they
admit infeasible plans. At this scale the claim that exact-state formulations
never violate the limit holds for every incumbent returned, but optimality
proofs need a longer limit than we used, and we do not draw runtime
conclusions from time-limited cells.

\subsection{Larger A-only formulation runs}
\label{subsec:large_a_only}

We also ran the four historical A-only formulations and the two uniform
flexible-start encodings on the same twelve instances with 20 or 30 aircraft
and 300--452 flights (three seeds per fleet-size/horizon cell), using Gurobi
and a 600\,s limit. This available run covers 12 instances, not the originally
planned 40-instance grid; P=40 and P=50 results below are separate 200\,s
scale checks and are not pooled with it. The first four formulations use the historical
zero-day-deferral setting; the uniform runs use their separate uniform
maintenance policy and are not directly comparable in objective value. The
original paper row and endpoint split were optimal in all twelve historical
runs, but every returned schedule failed the independent flight-hour replay.
The strengthened legacy and optimized event formulations were each optimal in
nine of twelve runs; all returned schedules passed replay. For the uniform
encodings, each proved optimality in only three of twelve runs. Among
time-limited uniform runs, mean reported gaps were 2.03\% for the day encoding
and 1.52\% for the event encoding. Replay passed for the available incumbents;
five day-encoding runs and six event-encoding runs had no schedule to replay.
These are scale checks, not reliable runtime rankings: several formulations
hit the time limit and the policies differ.

\begin{table}[t]
\centering\scriptsize
\caption{Gurobi results on the larger A-only instance set. ``Opt.'' is the
number of runs proven optimal out of 12; mean gaps are for time-limited runs
only. Replay reports the independent flight-hour audit.}
\label{tab:large_a_only}
\begin{tabular}{p{0.34\linewidth}rrrr}
\toprule
Formulation & Opt. & Mean wall (s) & Mean limit gap (\%) & Replay \\
\midrule
Legacy paper C13 & 12/12 & 20.2 & -- & Fail, 12/12 \\
Endpoint split & 12/12 & 37.8 & -- & Fail, 12/12 \\
Strengthened legacy & 9/12 & 94.4 & 0.90 & Pass, 12/12 \\
Optimized event & 9/12 & 175.7 & 0.41 & Pass, 12/12 \\
Uniform flexible day & 3/12 & 515.3 & 2.03 & Pass, 7/12; 5 not run \\
Uniform flexible event & 3/12 & 582.0 & 1.52 & Pass, 6/12; 6 not run \\
\bottomrule
\end{tabular}
\end{table}

\subsection{Discussion and limitations}
\label{subsec:robustness}

Three findings are supported by the experiments. First, C5 removes the
non-executable assignments addressed by the corner case, which occurs in the
timetables used. Its computational cost depends on instance size and costs:
larger instances have more rows and can require more solution time
(\Cref{subsec:c5_impact}). Second, schedules
produced by the published cumulative-flight-time row and its endpoint split in
isolation violate the flight-hour limit in every replayed instance, with both
solvers, whereas the exact-state formulations never do
(\Cref{subsec:solver_sensitivity}). Third, the day- and event-indexed uniform
encodings return identical optima and, on this family, comparable solution
times (\Cref{subsec:ten_seed}).

The study has clear limits. All instances are synthetic. The matched families
contain at most 240 flights and 15 aircraft; selected scale checks extend to
752 flights and 50 aircraft, but do not reproduce the full Khaled et al.\
benchmark protocol. The experiments use a single maintenance hub and A checks
only, and no real airline timetable has been evaluated. The six
matched cells use one seed each; the replication subsets cover only the
P=5/10 cells. CPLEX~12.6 and Gurobi~11.0.3 are used with one configuration
each, so timing comparisons are descriptive, and the replay counts depend on
which optimal schedule the solver returns. Larger fleets, varied airport
networks and station capacities, non-zero initial flight-hour counters, B/C/D
checks, and the full nightly and calendar-spacing policy of Khaled et al.\ are
not covered.

%%==========================================================================
\section{Conclusions and Future Work}
\label{sec:conclusion}
%%==========================================================================

This paper revisited the compact Tail Assignment Problem formulation of
\citet{khaled2018compact} and made three focused contributions. First, we
showed that the basic continuity and turn-time inequalities admit a
simultaneous-departure corner case and closed it with explicit non-overlap
constraints. Second, we analysed the missing boundary protections of using
the original conditional hour row alone and the failure of its endpoint-split
alternative. Exact states guard each pre-check arrival; strict day-boundary
opening rows now include the first day and flights through their ending day.
Third, we formulated uniform day-indexed and event-indexed encodings for
flexible-start, hour-limited A checks, with cross-day deferral and at most one
check per aircraft per actual start day. They share an event-aware hour
engine and are equivalent by direct substitution. This is an encoding result,
not independent physical validation or a day-free reformulation of Khaled's
complete nightly/calendar-spacing model. A fixed 22:00 comparison restricts
the flexible model without changing its resource or cost rules.

The practical significance of the first correction depends on the instance:
the timetables used here contain up to 153 same-airport pairs with identical
departure times, yet with aircraft-specific costs the legacy optimum was
already executable. Under neutral costs the legacy rows produced a cheaper but
non-executable plan in four of seven instances, with an objective understated
by 0.008--0.035\%. On twelve larger instances (300--452 flights, 20--30
aircraft) the legacy rows were non-executable under neutral costs in every
instance, with an objective understated by 0.04--0.14\%, whereas C1--C5 never
was. The non-overlap rows therefore cost model size and, at this scale, some
solution time, but they are needed for executable plans and change the
objective only slightly.

In the historical, nonuniform-policy study, post-solve replay reports no
flight-hour violations for the six
strengthened legacy optima and the four optimal schedules from each Event-based
method in the CPLEX run. The original paper-C13 and endpoint-split
formulations produce 663 and 183 flight-hour replay violations, respectively,
across the six CPLEX-selected schedules. A matched Gurobi~11.0.3 rerun returns
the same objective values but yields 672 and 200 violations, respectively.
These differing totals show that replay counts depend on the selected optimal
schedule and should not be interpreted as invariant formulation properties.
Nevertheless, every replayed schedule from these two formulations fails the
independent hour audit. All six corrected schedules and the four optimal
Event-based schedules per method pass in the CPLEX run; all six Gurobi
schedules per corrected or Event-based method pass. Gurobi also proves
optimality for the four larger Event-based cases that reached the CPLEX time
limit, and those Gurobi schedules pass replay. The audit checks event timing,
capacity, flight blocking, and flight-hour feasibility. New uniform runs
replay every loaded incumbent, including time-limited ones, and save the
decisions for a separate artifact-level cost and feasibility audit.

The larger A-only runs extend the scale check to 752 flights and 50 aircraft.
They reinforce the need for independent replay: the original row and endpoint
split fail the hour audit in every returned schedule, while all available
strengthened-legacy and event-based incumbents pass. The uniform encodings
prove optimality in only three of twelve runs each at 300--452 flights, so
their scale results are limited to feasibility of returned incumbents and do
not support a runtime ranking. No real airline timetable was available for
this study.

The historical results do not show a general Event-based speedup. On the four common
optimal cells, the strengthened legacy model has lower mean CPU and wall time
than either Event-based method, while the per-cell and model-size results vary.
An optional A-check-pruning reduction also preserves the optimum on one
constructed instance where no tail can reach its hour threshold, but this is
model-size evidence rather than a runtime result on maintenance-critical cases.
The six-cell scale-up uses one seed per cell; a ten-seed Gurobi replication of
the P=5/10 cells returns identical optima for the two uniform encodings in all
80 runs and finds the event-indexed encoding faster in 6--9 of ten instances
per cell, with modest effect sizes and a contrary sign in the earlier CPLEX
three-seed subset. This evidence is descriptive and does not
establish a statistical performance ranking or solver-independent dominance.
A matched Gurobi run supports the objective and feasibility conclusions, but
its solver-selected replay counts differ for the two deficient hour
formulations. The computational evidence uses one configuration per solver;
timings remain descriptive and do not establish solver-independent or
statistically isolated performance differences.

\paragraph{Reproducibility.}
The project contains the model implementations, deterministic instance
generator, benchmark scripts, result tables, source snapshots, saved schedules,
provenance sidecars and plotting scripts needed to reproduce the reported
snapshots. The legacy, event, and
ordered A-only builders remain separately selectable so comparisons can
distinguish model semantics from implementation choices.

\subsection*{Future directions}

\begin{itemize}
  \item \textbf{Larger A-only benchmark families:} repeat the horizon and
    fleet-size experiments over multiple seeds and heterogeneous timetable
    structures, reporting solver variability and confidence intervals.
  \item \textbf{Solver and configuration robustness:} extend the initial
    CPLEX--Gurobi cross-check across multiple solver configurations, including
    presolve, time limits, and warm-start policies, and report variability in
    both runtimes and replay outcomes.
  \item \textbf{Integrated fleet assignment and routing:} extend the corrected
    compact model while retaining physically executable route constraints.
  \item \textbf{Robust and recoverable variants:} study uncertainty in flight
    times, disruptions, and maintenance availability under the exact state
    formulation.
  \item \textbf{Heuristic methods:} develop constructive and local-search
    methods for the corrected A-only model and compare them with the exact
    formulations on larger instances.
\end{itemize}

%% ==========================================================================

\bibliographystyle{elsarticle-harv}
\bibliography{refs}

\end{document}
